\documentclass[11pt,a4paper]{article}
\usepackage[utf8]{inputenc}
\usepackage[T1]{fontenc}
\usepackage[ngerman]{babel}
\usepackage{amsmath,amsthm,amssymb}
\usepackage{mathtools}
\usepackage[margin=2.3cm]{geometry}
\usepackage{hyperref}
\usepackage{enumitem}
\usepackage{xcolor}
\usepackage{tikz}
\usepackage{booktabs}

% --- Theorem environments ---
\theoremstyle{plain}
\newtheorem{theorem}{Theorem}[section]
\newtheorem{proposition}[theorem]{Proposition}
\newtheorem{lemma}[theorem]{Lemma}
\newtheorem{corollary}[theorem]{Korollar}
\newtheorem{conjecture}[theorem]{Vermutung}
\newtheorem{blueprint}[theorem]{Bedingter Entwurf}
\newtheorem{milestone}[theorem]{Meilenstein}
\newtheorem{hypothesis}[theorem]{Hypothese}

\theoremstyle{definition}
\newtheorem{definition}[theorem]{Definition}
\newtheorem{remark}[theorem]{Bemerkung}
\newtheorem{example}[theorem]{Beispiel}

% --- Operators ---
\DeclareMathOperator{\spec}{spec}
\DeclareMathOperator{\Spec}{Spec}
\DeclareMathOperator{\Dom}{Dom}
\DeclareMathOperator{\Gal}{Gal}
\DeclareMathOperator{\Tr}{Tr}
\DeclareMathOperator{\Res}{Res}
\DeclareMathOperator{\vol}{vol}
\DeclareMathOperator{\id}{id}
\DeclareMathOperator{\sgn}{sgn}
\DeclareMathOperator{\Hom}{Hom}
\DeclareMathOperator{\End}{End}
\DeclareMathOperator{\Aut}{Aut}
\DeclareMathOperator{\Rep}{Rep}

% --- Shortcuts ---
\newcommand{\C}{\mathbb{C}}
\newcommand{\R}{\mathbb{R}}
\newcommand{\Q}{\mathbb{Q}}
\newcommand{\Z}{\mathbb{Z}}
\newcommand{\N}{\mathbb{N}}
\newcommand{\A}{\mathbb{A}}
\newcommand{\F}{\mathbb{F}}
\newcommand{\HH}{\mathcal{H}}
\newcommand{\V}{\mathcal{V}}
\renewcommand{\S}{\mathcal{S}}
\newcommand{\E}{\mathcal{E}}
\newcommand{\D}{\mathcal{D}}
\newcommand{\K}{\mathcal{K}}
\newcommand{\M}{\mathcal{M}}
\newcommand{\T}{\mathcal{T}}
\newcommand{\chifn}{\chi}
\newcommand{\gap}{\operatorname{gap}}
\renewcommand{\Re}{\operatorname{Re}}
\renewcommand{\Im}{\operatorname{Im}}

% --- Metadata ---
\title{Funktionale Stabilitätstheorie --- Physikalische Instanziierung durch das\\
Renormierte Freie-Energie-Prinzip\\
\vspace{0.3em}
{\large [RFEP v1.8 Update Entwurf --- Post-Zookeeper Transfer]}}
\author{Lukas Geiger\\
\small Unabhängiger Forscher, Bernau}
\date{26. April 2026}

\begin{document}

\maketitle

\begin{abstract}
Die Funktionale Stabilitätstheorie (FST) ist ein Programm, das ein einziges
strukturelles Muster --- funktionale Positivität unter einer Eichbedingung ---
als das gemeinsame Substrat offener Probleme in der Zahlentheorie, der mathematischen Physik
und der Kosmologie identifiziert. Dieses Paper ist die physikalische Instanziierung der FST.
Wir präsentieren das Renormierte Freie-Energie-Prinzip (RFEP) als die universelle Stabilitäts-Metaregel,
die dem Muster~A (PA1--PA3) und dem Dissipativen Selektionsprinzip (DS1--DS3) zugrunde liegt.

Es werden drei ontologische Ebenen unterschieden: (1) das abstrakte RFEP als
universelle Stabilitäts-Metaregel; (2) das Dissipative Selektionsprinzip
als Mechanismus, durch den die physikalische Evolution den eindeutigen
dissipationsminimierenden Attraktor auswählt; und (3) domänenspezifische
Instanziierungen als konkrete Verifikationen. Wir überprüfen das Framework
in fünf Bereichen: FST-Physik (Yang--Mills, Navier--Stokes, Navier--Stokes Log-Distanz-Ungleichung,
Turbulenz K41) und FST-Kosmologie (Dunkle Energie via das Kooperative Renormierungs-Modell).
Der historische Ursprung in der Riemannschen Vermutung und das mathematische Fundament
der FST werden im Begleitpapier \cite{ZetaZooMaster} entwickelt. Veröffentlicht in englischer und deutscher Fassung.
\end{abstract}

\tableofcontents

% =====================================================
\section{Einleitung und Problemstellung}
\label{sec:intro}
% =====================================================

\subsection{Die Spektrum-Nullstellen-Identifikation}

Hilbert und P\'olya schlugen vor, dass die Riemannsche Vermutung (RH) aus der
Existenz eines selbstadjungierten Operators $H$ auf einem Hilbertraum $\HH$ folgen würde,
dessen Spektrum die Menge der Imaginärteile der nicht-trivialen Nullstellen
der Riemannschen Zetafunktion ist:
\begin{equation}
  \spec(H) = \{\gamma_\rho \in \R : \zeta(1/2 + i\gamma_\rho) = 0\}.
\end{equation}
Für Dirichlet: finde $H_\chifn$ mit $\spec(H_\chifn) =
\{\gamma_\chifn^{(k)} : L(1/2 + i\gamma_\chifn^{(k)}, \chifn) = 0\}$.

Diese \emph{Spektrum-Nullstellen-Identifikation} ist der offene Schritt im
Connes-Programm \cite{Connes1999,Connes2026,ConnesConsaniMoscovici2025}.

\subsection{Zwei Hindernisse auf der direkten Hilbert-Seite}

\begin{enumerate}[leftmargin=*]
  \item \textbf{Eigenvektor-Analytizitätsbarriere.} Für den Dirichlet-Weil-Operator
    auf $L^2([-L,L])$ ist die Paley-Wiener-Fortsetzung durch $\sigma = 1/2$ begrenzt
    (archimedischer Pol, Prim-Dirichlet-Konvergenz) \cite{AnalyticKernelV2,C2LayerStatus}.
    Die L-Nullstellen treten auf dieser Ebene nicht in Erscheinung.
  \item \textbf{NE-B Gruppenstruktur-Hindernis.} Die Prim-Verschiebungs-Algebra
    auf $L^2([-L,L])$ hat eine Zentralisator-Dimension von $1$
    \cite{Geiger2026RHII}, Abschnitt~5. Die Halbgruppen-Gruppen-Äquivalenz-Vermutung
    (SGE; \cite{MathMaster}, Abschnitt~9) prognostiziert, dass der direkte
    Hilbert-P\'olya-Ansatz ohne eine Gruppenstruktur scheitert.
\end{enumerate}

Diese Kombination motiviert eine Nicht-Hilbert-, Nicht-direkt-Galerkin-Realisierung.
Meyer 2005 erreicht dies über eine Idele-Klassen-Darstellung auf einem
Fr\'echet-Distributionsraum. FST ist ein Vorschlag für ein axiomatisches Vokabular,
in dem Meyers Konstruktion die natürliche abelsche Instanziierung ist und in dem
alternative Routen (CCM-Prolat, Bost-Connes-Thermodynamik) verglichen werden können.

\subsection{Beitrag}

Dieses v1.8 Update reorganisiert das Paper um das RFEP als die physikalische
Instanziierung der Funktionalen Stabilitätstheorie. Sein Beitrag ist vierfach:
\begin{enumerate}[leftmargin=*]
  \item Es trennt die ontologischen Ebenen RFEP, DS1--DS3 und die domänenspezifische Instanziierung;
  \item Es dokumentiert die Verpflichtungen des Post-Zookeeper-Transfers als explizite Prüfungen
    von Operator/Form, Defekt, Galerkin-Treue, Komplement-Lücke und Spur-Kompatibilität;
  \item Es tabellarisiert den aktuellen Verifikationsstatus über die Zweige Physik und Kosmologie hinweg; und
  \item Es hält das ursprüngliche Spektrum-Nullstellen/RH-Material zweiglokal und
    verweist für das mathematische Fundament auf das Begleitpapier Zeta Zoo \cite{ZetaZooMaster}.
\end{enumerate}

Es wird kein neues Theorem über RH, Yang--Mills, Navier--Stokes oder Kosmologie behauptet.
Der Hauptbeitrag ist eine statusstabile RFEP-Ebene: sie unterscheidet bewiesene
strukturelle Reduktionen, bedingte Transfers, numerische Machbarkeitsnachweise
(Proof-of-Life) und offene analytische Verpflichtungen.

\subsection{Bezug zu früheren Versionen}

Dies ist der Entwurf des RFEP v1.8 Updates. Er integriert das v0.8 Post-Zookeeper-Transferpaket,
ändert jedoch die Rolle auf Ebene des Papers: RFEP wird nun als physikalische Instanziierung
der FST präsentiert, anstatt als Anhang zur Spektrum-Nullstellen-Routemap. Gegenüber v0.8
fügt das Update die Zweiglokalität, die Trennung des Zeta Zoos, die Status-Tabelle
für Physik/Kosmologie-Beweise und den CRM-Kreis im Ausblick hinzu. Gegenüber v0.7
behält es das Drei-Lemma-Endspiel und die RFEP-Transferverpflichtungen für physikalische
Nachträge bei. Gegenüber v0.6 bleiben Behauptungen über Theorem 8.4 (MS1 $\Leftrightarrow$ MS2),
Vermutung 8.6 (Weyl-Gesetz für $QW_\lambda$), Proposition 8.2 (CCM instanziiert (FST-7)(B)),
die Spurklasse-Behauptung für $G_\beta^{\mathrm{ren}}$ und die Rhetorik der Residuallücke
im Emergenz-Abschnitt herabgestuft, um der technischen Strenge zu entsprechen, die
tatsächlich in den unterstützenden Notizen vorhanden ist. Ein externes Audit (April 2026)
identifizierte systematische Übertreibungen in v0.6; die vorliegende Version hält
die Aussagen des Papers auf einer Linie mit den Diagnosen auf Notizenebene.

% =====================================================
\section{FST-Selbstdualität ohne Hilbertraum}
\label{sec:fst-duality}
% =====================================================

\subsection{Was ersetzt die Selbstadjungiertheit?}

Selbstadjungiertheit auf einem Hilbertraum erzwingt gleichzeitig ein reellwertiges Spektrum
und vollständige Eigenvektorfamilien. Auf einem Nicht-Hilbertraum müssen diese
separat erzwungen werden.

\begin{definition}[FST-Selbstdualität]
\label{def:fst-duality}
Ein Tripel $(\V, B, G)$ ist ein \emph{FST-selbstduales System}, wenn:
\begin{enumerate}[leftmargin=*]
  \item $\V$ ein lokalkonvexer topologischer Vektorraum ist (typischerweise
    Fr\'echet oder reflexiv Banach).
  \item $B: \V \times \V \to \C$ eine stetige nicht-ausgeartete
    Sesquilinearform (die \emph{Selbstdualitätspaarung}) ist.
  \item $G$ eine lokalkompakte Gruppe ist, die stetig auf $\V$ operiert
    und $B$ erhält.
\end{enumerate}
\end{definition}

Wenn $\V$ ein Hilbertraum und $B$ das Skalarprodukt ist, ist dies eine unitäre
$G$-Darstellung.

\begin{proposition}[Reelles Spektrum durch $B$-Symmetrie]
\label{prop:real-spectrum}
Sei $(\V, B, G)$ FST-selbstdual. Wenn $H$ der infinitesimale
Erzeuger einer einparametrigen Untergruppe ist, die $B$-symmetrisch operiert
($B(Hv, w) = B(v, Hw)$ auf einem dichten Definitionsbereich), dann sind die gemeinsamen
Eigenwerte von $H$ bezüglich $B$ reell.
\end{proposition}

Das Argument ist Standard für $B$-symmetrische Operatoren auf reflexiven
Räumen \cite{Meyer2005}.

\subsection{Warum FST Nicht-Hilbert motiviert}

FST-Stabilität (RFEP; \cite{MathMaster}, Abschnitt~6.5): Ein System ist stabil genau dann, wenn sein RG-Fluss
gegen einen selbstähnlichen Fixpunkt konvergiert.

Für eine Operatorfamilie $\{A_\lambda\}_{\lambda > 0}$ ist Stabilität
unter RG-Skalierung im Allgemeinen nicht mit einer Hilbert-Struktur kompatibel.
Ein Fr\'echet- oder lokalkonvexer Raum, der die Reskalierung als
Automorphismus zulässt, ist natürlich. Die Idele-Klassengruppe ist
das prototypische Beispiel.

% =====================================================
\section{Formalisierung der FST-Axiome}
\label{sec:axioms}
% =====================================================

\subsection{Grunddaten}

Eine \emph{FST-Struktur} besteht aus:

\begin{description}[style=nextline]
\item[(G)] Einer lokalkompakten abelschen topologischen Gruppe $G$ mit Haar-Maß.
\item[(R)] Einer abgeschlossenen Untergruppe $R \leq G$, isomorph zu $\R$ oder
  $\R_+^\times$ (die \emph{Renormierungsachse}).
\item[(N)] Einer komplementären abgeschlossenen Untergruppe $N \leq G$ mit $G
  \cong R \times N$, $N$ kompakt-proendlich (der \emph{arithmetische Nukleus}).
\item[($\chifn$)] Einem unitären Charakter $\chifn : G \to \C^\times$.
\item[(S)] Einer $G$-äquivarianten Bruhat-Schwartz-Klasse $\S(G)$, mit
  raschem Abfall entlang $R$, lokal konstanter Regularität entlang $N$,
  und einer Fourier-Transformation $\mathcal F : \S(G) \to \S(\hat G)$
  mit Pontryagin-Dualität $\hat G \cong G^\ast$.
\end{description}

\subsection{Muster-A Axiome}

\begin{description}[style=nextline]
\item[(P1)] $G = R \times N$ topologisch.
\item[(P2)] $\S(G) = \S(R) \otimes \S(N)$ topologisch.
\item[(P3)] $\mathcal F = \mathcal F_R \otimes \mathcal F_N$.
\end{description}

\subsection{Das SGE Axiom}

\begin{description}[style=nextline]
\item[(SGE)] $G$ ist eine Gruppe (keine Halbgruppe). Translationen $T_g:
  \S(G) \to \S(G)$ sind invertierbar.
\end{description}

Dies ist das $G$-Analogon zu dem, was auf $L^2([-L, L])$ unter
Prim-Verschiebungen scheitert (vgl.~NE-B, \cite{Geiger2026RHII}, Abschnitt~5).

\subsection{Die DS1--DS3 Axiome}

\begin{description}[style=nextline]
\item[(DS1) \textit{Dissipation.}] Die $R$-Operation
  $\{T_r\}$ ist $\S$-erhaltend und kontraktiv in einer gewichteten Halbnorm.
\item[(DS2) \textit{Selektion.}] $P_\chifn : \S(G) \to
  \S(G)^\chifn$ ist stetig und idempotent.
\item[(DS3) \textit{Reproduktion.}] Für $N = \prod_\ell' N_\ell$,
  $\S(G)^\chifn = \S(R) \otimes \bigotimes_\ell \S(N_\ell)^{\chifn_\ell}$.
\end{description}

\subsection{Das RFEP Axiom}

\begin{description}[style=nextline]
\item[(RFEP)] Es existiert ein Funktional der freien Energie
  $\mathcal F_\chifn : \S(G)^\chifn \to \R$, das konvex und koerziv
  unter der $R$-Operation ist und dessen kritische Punkte die
  Eigendistributionen des $R$-Erzeugers sind.
\end{description}

Diese sechs Axiome bilden die \emph{FST-Basis}.

\begin{remark}[Status des Axiom-Pakets]
\label{rem:axioms-status}
Als Klassifizierungssprache ist das Axiom-Paket nützlich: Es trennt sauber ein Basis-Framework vom Vollständigkeitsprinzip (FST-7) und von später eingeführten stärkeren Selektoren (FST-Free, FST-Min), die erst in Abschnitt~\ref{sec:emergence} vorgestellt werden. Als minimales Axiomensystem ist es jedoch \emph{noch nicht} überzeugend: Minimalität und Unabhängigkeit werden behauptet, aber nicht bewiesen. Einige Punkte (z.B. Muster~A) überschneiden sich erheblich mit den Grunddaten; RFEP wird als zentrales Prinzip behandelt, aber seine Wirkungskraft relativ zu den anderen Axiomen wird nicht scharf isoliert; und die späteren Erweiterungen (FST-Free, FST-Min) fungieren eher als \emph{Selektionsprinzipien, die arithmetischen Inhalt kodieren}, denn als Konsequenzen. Die angemessene Lesart ist daher: Die Axiome sind eine nützliche Karte dafür, wo Inhalte verortet sind, aber noch kein fundamentales Axiomensystem mit bewiesener Unabhängigkeit.
\end{remark}

\subsection{Kanonische L-Funktion}

Für $f \in \S(G)^\chifn$,
\[
  \tilde f(s) := \int_R f(r, e_N) \, r^s \, \frac{dr}{r},
\]
\[
  \Lambda_{G,\chifn}(s) := Z_\infty(s) \cdot \prod_\ell
    L_\ell(s, \chifn_\ell).
\]
Dies ist die Tate-These \cite{Tate1950} axiomatisiert.

% =====================================================
\section{Zweiglokalität und der Zeta Zoo}
\label{sec:branch-locality}
% =====================================================

RFEP v1.8 verwendet eine zweiglokale Lesart der FST. Der mathematische
Zweig wuchs historisch aus der Riemannschen Vermutung, den Hilbert--P\'olya
Programmen und dem Zeta Zoo. Dieser Zweig liefert das stärkste
formale Vokabular: Gruppenstruktur, Selbstdualität, Eichbedingungen,
Positivitätsformen und explizite Spur-Hindernisse. Er ist jedoch nicht dasselbe
Objekt wie die physikalische RFEP-Ebene.

In diesem Paper ist RFEP die physikalische Metaregel:
\[
  \text{stabile physikalische Evolution}
  \quad\Longleftrightarrow\quad
  \text{die renormierte freie Energie erreicht ein Minimum unter Nebenbedingungen.}
\]
Das Dissipative Selektionsprinzip (DS1--DS3) ist der Mechanismus, durch den
das RFEP-Minimum dynamisch ausgewählt wird. Domänenspezifische Papiere instanziieren diesen Mechanismus dann mit konkreten Funktionalen, Operatoren und Defektstrukturen.

Das Begleitpapier Zeta Zoo \cite{ZetaZooMaster} wird daher als das mathematische Fundament und der historische Ursprung der FST behandelt, nicht als Voraussetzungsbeweis für die hier gemachten physikalischen Behauptungen. Diese Trennung verhindert eine häufige Fehlinterpretation: Physikalische RFEP-Anwendungen können Vokabular und Beweisverpflichtungen aus dem Zeta Zoo erben, aber sie erben keine RH-Behauptung.

\subsection{Aktueller Domänenstatus}

\begin{table}[h]
\centering
\begin{tabular}{p{0.22\textwidth} p{0.24\textwidth} p{0.42\textwidth}}
\toprule
Domäne & RFEP-Status & Gegenwärtiger Beweisstatus \\
\midrule
Yang--Mills & Starke Kopplung RFEP Instanz &
Starke Kopplung Massenlücke auf dem Gitter in den aktuellen Notizen bewiesen; Kontinuums-Lücken-Transfer bleibt bedingt. \\
Navier--Stokes & Bedingte RFEP/Defekt Instanz &
$L^2$-Hindernistheorem und numerische Reichweiten-Diagnosen vorhanden; globale Regularität bleibt bedingt durch Annahme~G und Bedingung~D. \\
NS Log-Distanz Ungleichung & RFEP-kompatible LDI/TLL Brücke &
BV-Kettenregel bewiesen; PDE-Ebene TLL/LDI-Abschluss bleibt offen. \\
Turbulenz K41 & Variationelle RFEP Instanz &
K41-Minimierer-Theorem im Skelett unbedingt; anomale Dissipation und DFC/SP-Abschluss bleiben bedingt. \\
Dunkle Energie / CRM & Kosmologische RFEP Instanz &
Kooperatives Renormierungsmodell liefert den beabsichtigten Kosmologie-Zweig; Beobachtungs- und Modell-Identifikationsprüfungen bleiben offen. \\
\bottomrule
\end{tabular}
\caption{Zweiglokaler Beweisstatus für das RFEP v1.8 Update.}
\label{tab:domain-status}
\end{table}

% =====================================================
\section{Der FST Fixpunktsatz}
\label{sec:fixed-point}
% =====================================================

\subsection{Die Konstruktion}

\begin{description}[style=nextline]
\item[$\V_\chifn$] Quotient von $\S(G)^\chifn$ nach Koinvarianten $\K$
  (Funktionen mit verschwindenden Tate-Integralen auf der kritischen Geraden).
\item[$B_\chifn$] Die Tate-Paarung $B_\chifn(v, w) = \int_G v(g)
  \mathcal F[w](g) \chifn(g) dg$.
\item[$H_\chifn$] Erzeuger der $R$-Operation.
\end{description}

\subsection{Der Satz}

\begin{theorem}[FST Fixpunktsatz, bedingt durch (FST-7)]
\label{thm:fst-fixed-point}
Sei $(G, R, N, \chifn, \S)$ eine FST-Struktur, die die Basis-Axiome und das Vollständigkeitsaxiom (FST-7) aus Abschnitt~\ref{sec:fst7} erfüllt. Dann gilt:

\noindent(a) $\V_\chifn$ ist ein nicht-trivialer Fr\'echet-Raum.

\noindent(b) $B_\chifn$ ist nicht-ausgeartet und $R$-äquivariant.

\noindent(c) $H_\chifn$ ist $B_\chifn$-symmetrisch.

\noindent(d) $\spec_\mathrm{pt}(H_\chifn) = \{\gamma \in \R :
  \Lambda_{G,\chifn}(1/2 + i\gamma) = 0\}$.

\noindent(e) $v_\gamma \in \V_\chifn^\ast$ erfüllt
  $B_\chifn(v, v_\gamma) = \tilde v(1/2 + i\gamma)$.
\end{theorem}

\begin{corollary}[abstrakte RH, bedingt]
\label{cor:abstract-rh}
Unter den Voraussetzungen \emph{einschließlich (FST-7)} liegen die nicht-trivialen Nullstellen von $\Lambda_{G,\chifn}$ auf $\Re(s) = 1/2$.
\end{corollary}

Die Teile (a)--(c), (e) folgen im adelischen Fall aus der Tate-These.
Teil (d) --- die spektrale Identifikation --- erfordert (FST-7), und dies ist genau der Input, der nicht aus den Basis-Axiomen ableitbar ist. Die Aussage ist daher eine \emph{Umformulierung} des klassischen Meyer/Tate-Inhalts in die FST-Sprache, kein neues Theorem.

% =====================================================
\section{Minimalbeispiel: Riemannsche Zetafunktion}
\label{sec:minimal-example}
% =====================================================

Die minimale nicht-triviale Instanziierung: $G = \A_\Q^\times/\Q^\times$,
$\chifn = \chifn_0$, $R = \R_+^\times$, $N = \hat\Z^\times$.
Alle Axiome sind trivialerweise erfüllt.

\subsection{Tate Integrale}

Mit $f = f_\infty \otimes \prod_p f_p$, $f_\infty(x) = e^{-\pi x^2}$,
$f_p = \mathbf{1}_{\Z_p}$:
\[
  Z_\infty(s) = \pi^{-s/2} \Gamma(s/2), \qquad
  Z_p(s) = (1 - p^{-s})^{-1},
\]
\[
  Z(f, s) = \pi^{-s/2} \Gamma(s/2) \, \zeta(s) = \Lambda(s).
\]
Die adelische Selbstdualität $\hat f = f$ liefert direkt
$\Lambda(s) = \Lambda(1-s)$.

\subsection{Der Operator}

$H_{\chifn_0} = -x_\infty \partial_{x_\infty}$, Multiplikation mit
$-s$ auf der Mellin-Seite. Mellin-Auswertungen $v_\gamma: f \mapsto
Z(f, 1/2+i\gamma)$ liegen in $\V_{\chifn_0}^\ast$. Die Nicht-Trivialität
an Nullstellen ist (FST-7).
% =====================================================
\section{Das Vollständigkeitsaxiom (FST-7): Drei Sichtweisen}
\label{sec:fst7}
% =====================================================

\subsection{(FST-7) algebraisch: Pontryagin-Selbstdualität}

\begin{definition}[abstrakte Meyer-Bedingung]
\label{def:fst7-a}
(FST-7A): es existiert eine diskrete Untergruppe $\Xi \leq G$ mit:
\textit{(i)} $\Xi^\perp = \Xi$ unter der Pontryagin-Paarung;
\textit{(ii)} $G / (R \cdot \Xi)$ kompakt (Produktformel);
\textit{(iii)} $B_\chifn$ ist $\Xi$-invariant.
\end{definition}

Im adelischen Fall erfüllt $\Xi = \Q^\times$ (i)--(iii) mit
$\prod_v |x|_v = 1$ als Produktformel.

\subsection{(FST-7) informationstechnisch: Nyquist-Shannon}

\begin{definition}[informationstheoretische Selbstdualität]
\label{def:fst7-b}
(FST-7B): für alle $f \in \S(G)$ mit $\hat f \in \S(\hat G)$,
\[
  \sum_{\xi \in \Xi} |f(\xi)|^2 = \sum_{\eta \in \hat\Xi}
    |\hat f(\eta)|^2.
\]
\end{definition}

\begin{proposition}
\label{prop:B-iff-A}
(FST-7B) $\Leftrightarrow$ (FST-7A)(i) unter Standard-Pontryagin-Bedingungen an $(G, \Xi)$.
\end{proposition}

Interpretation: $\Xi$ ist die Nyquist-optimale Abtastung von $G$.

\begin{remark}[Slepian-Pollak-Prolate als natürliche Realisierung]
\label{rem:b-view-prolate}
Das Axiom (FST-7B) ist kompatibel mit dynamischen Realisierungen durch
Prolat-Sphäroid-Wellenfunktionen $h_{n,\lambda}$ von Slepian-Pollak
\cite{SlepianPollack1961}: diese minimieren gleichzeitig
Energie außerhalb des Intervalls und außerhalb des Bandes auf $\R$ und sind
daher eine Nyquist-Shannon-optimale Basis; für $n \equiv 0 \pmod 4$ sind sie
Eigenfunktionen der Fourier-Transformation, tragen also eine konkrete
Zeit-Frequenz-Selbstdualität. Dies ist der konzeptionelle Einstiegspunkt
für das CCM-Programm (vgl. Abschnitt~\ref{sec:ccm-prolate}). Die
\emph{Beziehung} zwischen dem Prolat-Setup und (FST-7)(B) ist eine
mathematisch sinnvolle Analogie und keine bewiesene Instanziierung;
siehe Bemerkung~\ref{rem:ccm-closest-model}.
\end{remark}

\subsection{(FST-7) dynamisch: arithmetische Ergodizität}

\begin{definition}[arithmetische Ergodizität]
\label{def:fst7-c}
(FST-7C): die $\Xi$-Operation auf dem Norm-$1$-Kern $G^{(1)}$ ist streng
ergodisch (Haar-Maß).
\end{definition}

\begin{proposition}
\label{prop:C-implies-A}
(FST-7C) mit Pontryagin-Fourier-invariantem Haar-Maß impliziert
die Poissonsche Summenformel für $\Xi$, also (FST-7A).
\end{proposition}

Im adelischen Fall operiert $\Q^\times$ durch strenge Approximation streng ergodisch auf
$\A^{(1)}$ \cite{CasselsFrohlich}.

\subsection{Äquivalenz als strukturelle Bedingungen}

\begin{theorem}[Äquivalenz der (FST-7) Sichtweisen]
\label{thm:fst7-equiv}
(FST-7A), (FST-7B), (FST-7C) sind äquivalent als strukturelle
Bedingungen für dasselbe $\Xi \leq G$ unter Standard-Pontryagin-
und Haar-Bedingungen.
\end{theorem}

\begin{remark}[Reichweite von Theorem~\ref{thm:fst7-equiv}]
Die Äquivalenz besteht auf der Ebene struktureller Eingangsbedingungen, nicht
auf der Ebene technischer Einstiegspunkte. In konkreten Fällen öffnen die drei
Sichtweisen sehr unterschiedliche technische Türen: (A) öffnet die Tate/Meyer-adelische Analyse;
(B) öffnet die Slepian-Pollack / Prolat-Operator-Analyse;
(C) öffnet BC-Thermodynamik / Cornelissen-Marcolli-Starrheit.
Die drei Einstiegspunkte sind als Beweisstrategien \emph{nicht} austauschbar.
\end{remark}

\subsection{(FST-7) ist aus Basis-Axiomen nicht ableitbar}

\textbf{Strukturelle Beobachtung.} (FST-7) in jeder seiner drei
Formen ist nicht allein aus (P1--P3), (SGE), (DS1--DS3), (RFEP) ableitbar.
Es ist eine arithmetische Eigenschaft, die von außen eingebracht werden muss:
aus der Tate-These im adelischen Fall, durch analogen Input in anderen Fällen.
Dies ist die zentrale ehrliche Einschränkung des Axiom-Pakets.

% =====================================================
\section{Die Bost-Connes/Meyer-Brücke}
\label{sec:bc-meyer}
% =====================================================

\subsection{Das BC95 System}

Die Hecke-$C^\ast$-Algebra $\HH_\Q = C^\ast(\N^\times \rtimes \hat\Z^\times)$
mit Zeitentwicklung $\alpha_t(\mu_n) = n^{it}\mu_n$ besitzt:
\begin{itemize}
\item $\spec(H_{BC}) = \{\log n : n \in \N^\times\}$.
\item $Z_{BC}(\beta) = \zeta(\beta)$ für $\beta > 1$.
\item KMS-Zustände: eindeutig für $\beta \leq 1$, parametrisiert durch
  $\hat\Z^\times \cong \Gal(\Q^\mathrm{ab}/\Q)$ für $\beta > 1$.
\item Phasenübergang bei $\beta = 1$.
\end{itemize}

\subsection{Vier Brückenansätze, einer empfohlen}

\begin{table}[h]
\centering
\begin{tabular}{@{}lccc@{}}
\toprule
\textbf{Ansatz} & \textbf{Plausibilität} & \textbf{Kurzfristig} & \textbf{Langfristig} \\
\midrule
1. GNS + Niedrigtemperatur-Limes & mittel & mittel & hoch \\
2. \textbf{Laplace-Mellin-Dualität} & \textbf{hoch} & \textbf{hoch} & mittel \\
3. Arithmetischer Site als Intermediär & hoch & niedrig & sehr hoch \\
4. Tropikalisierung & niedrig & niedrig & mittel \\
\bottomrule
\end{tabular}
\caption{Brückenansätze. Ansatz~2 ist das empfohlene kurzfristige Ziel;
er impliziert \emph{keinen} Beweis der RH, bietet aber eine strukturelle Identifikation
zwischen der thermodynamischen und der distributionellen Seite.}
\label{tab:bridge-approaches}
\end{table}

\subsection{Laplace-Mellin-Dualität: die Identifikation}

\begin{equation}
\label{eq:log-z-deriv}
  L(\beta) := -\frac{d}{d\beta} \log Z_{BC}(\beta)
    = -\frac{\zeta'(\beta)}{\zeta(\beta)}.
\end{equation}
Über die explizite Formel ist $L$ meromorph mit:
\begin{itemize}
  \item Einfachem Pol bei $\beta = 1$ (Mangoldt-Pol), Residuum $+1$.
  \item Polen bei $\beta = \rho$ (nicht-triviale Nullstellen), Residuen $-m_\rho$.
  \item Polen bei trivialen Nullstellen $\beta = -2k$.
\end{itemize}

\begin{conjecture}[Meyer-Eigendistributionen als BC-Residuen]
\label{conj:meyer-bc}
Für die Riemannsche Zetafunktion ($\chifn = \chifn_0$),
\[
  v_{\gamma_\rho}(f) = \Res_{\beta = 1/2 + i\gamma_\rho}
    L(\beta) \cdot \mathcal M[f](\beta).
\]
\end{conjecture}

Dies ist eine konjekturale Identifikation, keine bewiesene Äquivalenz.
Sie rigoros zu machen, ist exakt der Inhalt der Wick-Mellin-Hindernisse
(Abschnitt~\ref{sec:wick-mellin-obstacles}); siehe dort.

% =====================================================
\section{Das CCM-Prolat-Modell: Die nächstgelegene bekannte Realisierung von (FST-7)(B)}
\label{sec:ccm-prolate}
% =====================================================

Das Connes-Consani-Moscovici (CCM)-Programm
\cite{ConnesConsani2023,ConnesMoscovici2022,ConnesConsaniMoscovici2025}
realisiert eine Störung vom Rang eins eines Skalierungsoperators,
dessen Spektrum numerisch gegen die nicht-trivialen Nullstellen von $\zeta$
konvergiert. Wir fassen die Konstruktion zusammen und diskutieren ihre
Beziehung zu (FST-7)(B); danach präsentieren wir den Entwurf der Route~B
für die Sprache der fehlenden CCM-Schritte, wobei MS2/C2 als aus dem
Zookeeper-Abschluss importiert markiert wird und MS1 als das
verbleibende Element in dieser älteren Routen-Sprache verbleibt.

Wir betonen vorab: Bis auf das unten importierte Zookeeper MS2/C2-Ergebnis
ist alles ein Entwurf und/oder eine Analogie.

\subsection{Die Konstruktion}
\label{sec:ccm-construction}

Fixiere $\lambda > 1$ und setze $L = 2 \log\lambda$.

\begin{description}[leftmargin=*,style=nextline]
\item[\textbf{Skalierungs-Tripel.}] Auf $L^2([\lambda^{-1}, \lambda],
  d^\ast u)$ mit $d^\ast u = du/u$ besitzt der Skalierungsoperator
  $D^{(\lambda)}_{\log} = -i\,d/d(\log u)$ mit periodischen Randbedingungen
  die Eigenfunktionen $V_n(u) = L^{-1/2}\exp(2\pi i n \log(\lambda u)/L)$, $n \in \Z$.
\item[\textbf{Weil-Quadratische Form.}] Aus Weils expliziter Formel
  wird die Sesquilinearform $QW(f,g) = \Psi(f^\ast * g)$ mit $\Psi =
  W_{0,2} - W_\R - \sum_p W_p$ konstruiert. Eingeschränkt auf
  $L^2([\lambda^{-1}, \lambda], d^\ast u)$ tragen nur Primzahlen $p \leq
  \lambda^2$ bei.
\item[\textbf{Test-Unterraum $E_N$.}] Aufspann der $V_n$ mit $|n|
  \leq N$; Dimension $2N+1$.
\item[\textbf{Störung vom Rang eins.}] Sei $\epsilon_N$ der
  kleinste Eigenwert der Einschränkung $QW_\lambda^N$ und $\xi$
  der entsprechende Eigenvektor (angenommen als einfach und gerade unter $u
  \mapsto u^{-1}$). Sei $\delta_N \in E_N$ die Darstellung der Dirichlet-Kern-Auswertung
  am Rand, normiert durch $\delta_N(\xi) = 1$. Definiere
  \[
     D^{(\lambda,N)}_{\log} := D^{(\lambda)}_{\log}
       - |D^{(\lambda)}_{\log}\xi\rangle\langle\delta_N|.
  \]
\end{description}

\begin{theorem}[CCM 2025, Theorem 1.1]\label{thm:ccm-mainthm}
(i) $D^{(\lambda,N)}_{\log}$ ist selbstadjungiert auf $E'_N \oplus
E_N^\perp$ mit $E'_N = E_N/\C\xi$.
(ii) $\det_{\mathrm{reg}}(D^{(\lambda,N)}_{\log} - z) = -i
\lambda^{-iz}\widehat{\widehat\xi}(z)$.
(iii) $\widehat{\widehat\xi}(z)$ ist ganzzahlig, alle Nullstellen liegen auf der
reellen Achse und stimmen mit $\mathrm{spec}(D^{(\lambda,N)}_{\log})$ überein.
\end{theorem}

\begin{remark}[Numerische Evidenz]\label{rem:ccm-numerics}
Für $\lambda = \sqrt{14}$, $N = 120$, unter Verwendung nur der Primzahlen
$p \leq 13$, werden die ersten fünfzig nicht-trivialen Nullstellen von $\zeta(1/2+is)$
mit Fehlern im Bereich von $10^{-60}$ bis $10^{-6}$ rekonstruiert
(vgl.~\cite[Table~1]{ConnesConsaniMoscovici2025}). Die statistische
Unwahrscheinlichkeit einer solchen Übereinstimmung wird von CCM auf
$10^{-1235}$ geschätzt. Numerische Daten dieser Stärke sind \emph{zwingende Evidenz};
sie stellen kein Theorem dar, das RH reduziert.
\end{remark}

\subsection{Die $QW$-$PW$-Dualität}
\label{sec:ccm-duality}

CCM identifizieren eine strukturelle Dualität zwischen der Weilschen
Quadratischen Form $QW_\lambda$ und dem Prolat-Wellenoperator
\[
  PW_\lambda := -\partial_x\big((\lambda^2-x^2)\partial_x\big) +
    (2\pi\lambda x)^2,
\]
dessen Eigenfunktionen $h_{n,\lambda}$ die Prolat-Sphäroid-Wellenfunktionen
von Slepian-Pollak \cite{SlepianPollack1961} sind.

\begin{center}
\small
\begin{tabular}{@{}lll@{}}
\toprule
& \textbf{Weil $QW_\lambda$} & \textbf{Prolat $PW_\lambda$} \\
\midrule
Ursprung        & Zahlentheorie (Weil) & Informationstheorie (Shannon) \\
Regime          & Infrarot (niedrige Eigenwerte) & Ultraviolett (hohe Eigenwerte)\\
Zeta-Abdruck    & Niedrige Nullstellen von $\zeta$ & UV-Verhalten der Nullstellen\\
Optimalität     & Explizite-Formel-Positivität & Zeit-Frequenz-optimal \\
Referenz        & \cite{ConnesConsaniMoscovici2025} & \cite{ConnesMoscovici2022} \\
\bottomrule
\end{tabular}
\end{center}

\subsection{CCM als nächstgelegenes Modell für (FST-7)(B)}
\label{sec:ccm-as-fst7b}

Erinnerung an (FST-7)(B) aus Abschnitt~\ref{sec:fst7}: Es existiert eine
\emph{Nyquist-Shannon-optimale} diskrete Teilmenge $\Xi \le G$, auf der
eine Plancherel-artige Identität gilt. Die Prolat-Sphäroid-Wellenfunktionen
von Slepian-Pollak \emph{sind} die kanonische Nyquist-Shannon-optimale
Basis in der Zeit-Frequenz-Ebene.

\begin{remark}[CCM als nächstgelegenes Modell für (FST-7)(B)]\label{rem:ccm-closest-model}
Die CCM-Konstruktion $(D^{(\lambda,N)}_{\log}, QW_\lambda, PW_\lambda)$
ist das \emph{nächstgelegene bekannte konkrete Modell} von (FST-7)(B) in dem Sinn,
dass (a) es um die kanonische Nyquist-Shannon-optimale Slepian-Pollak-Basis
in der Zeit-Frequenz-Ebene herum aufgebaut ist, und (b) die Skalierungsoperation
von $\R_+^\ast$ auf $L^2([\lambda^{-1},\lambda], d^\ast u)$ das lokalkompakte abelsche
Umfeld liefert, das für die Formulierung von (FST-7)(B) erforderlich ist. Diese
Beziehung ist eine mathematisch sinnvolle Analogie, die die abstrakte
Optimalitätsaussage mit einem konkreten operatortheoretischen Rahmen in Einklang
bringt; sie ist \emph{keine} bewiesene Äquivalenz auf Theorem-Ebene. Eine
rigorose Aussage würde erfordern, ein explizites $\Xi$ aus Prolat-Daten zu
konstruieren, eine Plancherel-artige Identität im Sinne von Definition~\ref{def:fst7-b}
zu beweisen und dies mit der Rang-eins-Störungsbedingung an $\delta_N$ in
Verbindung zu bringen. Nichts davon ist hier etabliert.
\end{remark}

\begin{conjecture}[$QW\!\leftrightarrow\! PW$ als (FST-7)(A)$\leftrightarrow$(FST-7)(B) Äquivalenzoperator, strukturelle Form]\label{conj:qw-pw-fst7}
Es wird vermutet, dass die CCM-Dualität $QW_\lambda \leftrightarrow PW_\lambda$
zwischen Infrarot- und Ultraviolett-Regime eine dynamische Repräsentation
der abstrakten Äquivalenz $(\text{FST-7})(A) \Leftrightarrow (\text{FST-7})(B)$
aus Theorem~\ref{thm:fst7-equiv} ist. Konkret: Es sollte eine Transformation
existieren, die $QW_\lambda$-Eigenfunktionen im Limes $\lambda \to \infty$ auf
$PW_\lambda$-Eigenfunktionen abbildet und die die Pontryagin-Selbstdualität
von $\Xi = \Q$ innerhalb von $\A_\Q$ auf der Ebene spektraler Maße realisiert.
Das vorliegende Paper liefert einen Kandidaten für diese Transformation
(Abschnitt~\ref{sec:fst-transfer}), beweist aber nicht die Verschachtelungs-Eigenschaft (intertwining).
\end{conjecture}

\subsection{Die fehlenden CCM-Schritte}
\label{sec:ccm-missing-steps}

Zwei verbleibende technische Lücken wurden im CCM-Programm identifiziert:
\begin{description}[leftmargin=*,style=nextline]
\item[(MS1)] \textbf{Einfach-gerade} am kleinsten Eigenwert von $QW_\lambda$.
  Klassischerweise bekannt für den Prolat-Operator $PW_\lambda$ \cite{SlepianPollack1961};
  numerisch bestätigt für $QW_\lambda$; noch nicht bewiesen. Bleibt offen.
\item[(MS2)] \textbf{Approximationsqualität.} Der Poisson-Quasimode
  $k_\lambda$ muss die wahre $QW_\lambda$-Eigenfunktion $\xi_\lambda$
  gut genug approximieren, um spektrale Konvergenz zu übertragen.
  \textbf{Gelöst im Zookeeper Paper} via des Drei-Lemma-Mikrocluster-Abschlusses:
  Skalare säkulare Aufhebung, projizierte Poisson-Quasimode-Kontrolle
  und ein koerzives Komplement ergeben die Massenkonzentrations-Schranke
  $\|(I-P_{V_\lambda})k_\lambda\| \leq r_\lambda/g_* \to 0$,
  was zeigt, dass $k_\lambda$ mit beliebiger Präzision im resonanten Eigenraum lebt.
\end{description}

(MS1) bleibt ein offenes Problem in dieser älteren Routen-Sprache. (MS2)
wird im CORE nicht länger als offen behandelt: das Zookeeper Paper dokumentiert
den theorematischen Abschluss des Approximationsschritts im CCM-Fourier-Modell,
wobei der Standard-CCM-Transfer und der Hurwitz-Übergang dort behandelt werden.

\subsection{FST-Transfer: ein bedingter Reduktionsentwurf}
\label{sec:fst-transfer}

Wir konkretisieren nun Vermutung~\ref{conj:qw-pw-fst7} weit genug, um
einen \emph{bedingten Reduktionsentwurf} für die zwei CCM-Lücken zu formulieren.
Dieser Unterabschnitt ist \emph{kein} Theorem und beweist \emph{keine} Reduktion;
er ist ein Schema, in dem jeder Schritt offen und klar identifiziert ist.

\subsubsection{Zwei Kandidaten für Isometrien}

Sei $H_{PW} = L^2([-\lambda, \lambda], dx)$ mit Reflexion
$\rho(g)(x) = g(-x)$ und $H_{QW} = L^2([\lambda^{-1}, \lambda], d^\ast u)$
mit Inversion $\iota(f)(u) = f(u^{-1})$. Die Logarithmus-plus-Reskalierungs-Abbildung
\[
  \Phi_\lambda := \log^\ast \circ \, \sigma_\lambda :
    L^2([-\lambda, \lambda], dx) \xrightarrow{\sigma_\lambda}
    L^2([-\log\lambda, \log\lambda], dx)
    \xrightarrow{\log^\ast}
    L^2([\lambda^{-1}, \lambda], d^\ast u)
\]
ist eine \emph{kinematische} Isometrie, die $\rho$ und $\iota$ miteinander verknüpft ---
eine direkte Rechnung.

Aufbauend auf der expliziten CCM-Konstruktion \emph{schlagen wir als
dynamische Abbildung vor}
\[
  \Psi_\lambda : H_{PW} \longrightarrow H_{QW}, \qquad
    g \mapsto \big(\mathcal E(g)\big)\big|_{[\lambda^{-1},\lambda]},
  \qquad \mathcal E(g)(u) := u^{1/2} \sum_{n \ge 1} g(nu),
\]
d.h.\ Einschränkung der Epstein-Mellin-Faltung auf $[\lambda^{-1},\lambda]$.
Ob $\Psi_\lambda$ als Operator von $H_{PW}$ nach $H_{QW}$ beschränkt ist,
ob es ein dichtes Bild hat und was seine genauen adjungierten und
Verknüpfungs-Eigenschaften sind, ist selbst eine nicht-triviale Frage.
Wir sammeln dies als den ersten offenen Punkt:

\begin{milestone}[Definitions- und Formentheorie für $\Psi_\lambda$]
\label{ms:psi-domain}
Etabliere, dass $\Psi_\lambda$ ein wohldefinierter, dicht definierter
Operator $H_{PW} \to H_{QW}$ ist, und bestimme seinen Form-Definitionsbereich,
seine Beschränktheitsklasse und seine Verknüpfungsstruktur in Bezug auf
die Involutionen $(\rho, \iota)$. Beweise oder widerlege, dass
$\Psi_\lambda^\ast \Psi_\lambda$ eine kompakte Resolvente hat.
\end{milestone}

\subsubsection{Bedingter Reduktionsentwurf}

\begin{blueprint}[MS1 $\Leftarrow$ (MS2 $\wedge$ beschränkte Verknüpfung)]
\label{bp:fst-transfer}
Angenommen, Meilenstein~\ref{ms:psi-domain} wurde so gelöst, dass $\Psi_\lambda$
eine beschränkte Verknüpfung auf geeigneten Form-Definitionsbereichen ist.
Angenommen weiterhin, es existieren Konstanten $\mu_\lambda > 0$
und $\varepsilon_\lambda \to 0$ für $\lambda \to \infty$, so dass
\begin{equation}\label{eq:dynamic-isometry}
  \big\| QW_\lambda \circ \Psi_\lambda
    - \mu_\lambda \, \Psi_\lambda \circ PW_\lambda \big\|
  \;\le\; \varepsilon_\lambda
\end{equation}
als Operatoren auf dem geeigneten dichten Definitionsbereich (vgl.~(MS2)) gilt.
Dann ist das Folgende eine bedingte Schlussfolgerung, kein Theorem:
Nach Slepian-Pollak \cite{SlepianPollack1961} ist der kleinste Eigenwert
$\mu_{0,\lambda}$ von $PW_\lambda$ einfach mit gerader Eigenfunktion
$h_{0,\lambda}$; unter \eqref{eq:dynamic-isometry} und der
Verknüpfungshypothese würde ein Störungsargument vom Kato-Typ \cite{Kato}
einen einfachen $QW_\lambda$-Eigenwert $\epsilon_\lambda$ mit einem Eigenvektor
$\xi_\lambda$ liefern, der $\|\xi_\lambda - c \cdot \Psi_\lambda h_{0,\lambda}\| \le
C \varepsilon_\lambda$ erfüllt, und die Erhaltung der $\rho/\iota$-Graduierung würde
die Eigenschaft ``gerade'' liefern. In diesem Szenario würde (MS1) immer dann
folgen, wenn (MS2) gilt.
\end{blueprint}

\begin{remark}[Status des Entwurfs~\ref{bp:fst-transfer}]
\label{rem:blueprint-caveats}
Die bedingte Schlussfolgerung in Entwurf~\ref{bp:fst-transfer} beruht
auf drei offenen Punkten: (1) $\Psi_\lambda$ ist tatsächlich eine beschränkte
Verknüpfung (Meilenstein~\ref{ms:psi-domain}); (2) die Operator-Norm-Schranke
\eqref{eq:dynamic-isometry} gilt mit dem behaupteten Abklingen; (3)
das Kato-Störungsargument ist gleichmäßig auf den relevanten Form-Definitionsbereichen
anwendbar. Jeder Punkt ist nicht-trivial. Insbesondere erfordert die
Standardlesart des Satzes von Kato Operator-/Form-Konvergenz in einem
Norm- oder Form-Resolventen-Sinn, die aus \eqref{eq:dynamic-isometry} \emph{nicht}
automatisch hervorgeht. Der Entwurf ist am besten als Formulierung dessen
zu verstehen, was ein Beweis von MS1 via MS2 etablieren müsste, nicht als
ein Reduktionstheorem.
\end{remark}

\begin{remark}[Rolle der FST im Entwurf]\label{rem:fst-contribution}
Selbst in seiner bedingten Form illustriert Entwurf~\ref{bp:fst-transfer}
die Art und Weise, wie die FST beitragen soll: die Axiomäquivalenz
(FST-7)(A)$\leftrightarrow$(FST-7)(B) motiviert strukturell die Suche nach
einem Verknüpfungsoperator $\Psi_\lambda$ zwischen den beiden konkreten
Operatoren auf der CCM-Seite; die CCM-Konstruktion selbst liefert
die Kandidaten-explizite Form von $\Psi_\lambda$. Was fehlt, ist der
analytische Inhalt (Beschränktheit, Kontrolle der Operator-Norm, Anwendung der
Störungstheorie).
\end{remark}

\subsection{Meilenstein-Agenda für Route~B}
\label{sec:route-b-milestones}

Indem wir die heuristische Vermutung (``Weyl-Gesetz für $QW_\lambda$'') von
v0.6 durch eine strukturierte Meilenstein-Agenda ersetzen, zerfällt Route~B
in vier konkrete, aufeinanderfolgende Meilensteine im Sinne der Audit-Empfehlung:

\begin{milestone}[Definitions- / Formentheorie für $\Psi_\lambda$]
\label{ms:1}
(= Meilenstein~\ref{ms:psi-domain}.) Etabliere, dass $\Psi_\lambda$ ein
wohldefinierter, beschränkt invertierbarer Verknüpfungsoperator auf geeigneten
Form-Definitionsbereichen ist. Erster wirklich rigoroser Schritt der Route~B.
\end{milestone}

\begin{milestone}[Vergleich niedriger Eigenwerte via Min-Max]
\label{ms:2}
Beweise einen Min-Max-Vergleich zwischen den ersten endlich vielen
Eigenwerten von $QW_\lambda$ und $PW_\lambda$, \emph{ohne}
eine vollständige Weyl-Asymptotik zu verlangen. Dies ist das operatortheoretische
Minimum, das benötigt wird, um (MS1) gleichmäßig anzugreifen.
\end{milestone}

\begin{milestone}[Endlich-dimensionale CCM-Trunkierung]
\label{ms:3}
Beweise ein endlich-dimensionales CCM-Trunkierungstheorem innerhalb von $E_N$:
Kontrolliere die Einschränkung von $QW_\lambda$ auf $E_N$ und ihren Vergleich
zur Prolat-Einschränkung unter Verwendung von Standard-endlichdimensionaler
Störungstheorie (Katos Theorem ist dort sauber anwendbar). Dies ist der Ort, an dem
die CCM-Numerik tatsächlich lebt.
\end{milestone}

\begin{milestone}[Asymptotisches Regime]
\label{ms:4}
Erst nachdem Meilensteine~\ref{ms:1}--\ref{ms:3} etabliert sind:
Untersuche das asymptotische Verhalten von $N_{QW_\lambda}(\mu)$ für
$\mu \to \infty$ und $\lambda \to \infty$. Dies ist das Analogon zu
dem, was v0.6 spekulativ als ``Weyl-Gesetz für die Weil-Form'' bezeichnete,
und wird hier ausdrücklich als \emph{letzter} Meilenstein in der Sequenz
markiert, nicht als erster.
\end{milestone}

\begin{remark}[Warum diese Reihenfolge]
\label{rem:milestone-ordering}
Die Meilensteine \ref{ms:1}--\ref{ms:3} haben endlich-dimensionalen oder
form-theoretischen Charakter und sind mit existierender Operator-Theorie-Maschinerie
wirklich handhabbar. Meilenstein~\ref{ms:4} betrifft echte globale Asymptotiken
einer nicht-standardmäßigen quadratischen Form und wird viel weniger
gut verstanden. Der v0.6-Entwurf behandelte Meilenstein~\ref{ms:4}, als
wäre es fast ein Theorem; die vorliegende Revision ordnet die Agenda
ausdrücklich so um, dass die handhabbare Arbeit zuerst und das tiefe asymptotische
Element zuletzt kommt. Siehe die externen Audit-Notizen für die Motivation dieser
Umstrukturierung.
\end{remark}

% =====================================================
\subsection{Das Drei-Lemma-Endspiel für Route~B}
\label{sec:three-lemma-endgame}
% =====================================================

Wir präsentieren nun eine bedingte Auflösung von (MS2) via einer
endlich-dimensionalen Spektralanalyse des CCM Rang-eins-Operators.
Das Argument erfordert drei unabhängige Lemmas und liefert (MS2) als
Korollar durch eine Standard-Quasimode-zu-Projektor-Ungleichung.

\subsubsection{Setup und Notation}

Fixiere $\lambda > 0$. Sei $L = 2\log\lambda$, $N = \max(40,
\lceil 12L \rceil)$, und sei $E_N \subset L^2[0,L]$ der Fourier--Galerkin-Unterraum,
der von $\{e_n(x) = \sqrt{2/L} \sin(n\pi x/L)\}_{n=1}^N$ aufgespannt wird. Innerhalb
von $E_N$ nimmt der CCM-Operator die Form an:
\begin{equation}\label{eq:rank-one-structure}
  A^{(N)} \;=\; H^{(N)} + \tilde{C}\,|\tilde{u}\rangle\langle\tilde{u}|,
\end{equation}
wobei $H^{(N)}$ die Galerkin-Einschränkung des Kinetik-plus-Potential-Teils ist
und $\tilde{u} \in E_N$ die Poisson-Kern-Interaktion kodiert.
Der Poisson-Quasimode $k = k_\lambda^{(N)} \in E_N$ ist die normalisierte
Galerkin-Projektion des Poisson-Kerns $K(e^{x/\lambda})$.
Wir bezeichnen mit $u_0$ den kleinsten Eigenwert von $A^{(N)}$ und
mit $\{(u_j, q_j)\}_{j=0}^{N-1}$ das vollständige Eigensystem.

\begin{hypothesis}[Galerkin-Treue]
\label{hyp:galerkin}
Für jedes $\lambda$ repräsentiert das endlich-dimensionale Galerkin-Modell
$A^{(N(\lambda))}$ mit $N = \max(40, \lceil 12L \rceil)$ treu die Struktur der
spektralen Cluster des CCM-Operators in folgendem Sinn: Die Eigenwert-Reihenfolge,
die Rang-eins Säkulargleichung und die Cauchy-Verflechtung mit dem Hintergrundoperator
$H^{(N)}$ gelten mit Fehlern, die durch den Galerkin-Trunkierungsfehler
$\varepsilon_{\mathrm{Gal}} \sim e^{-c \cdot N/L}$ kontrolliert werden.
\end{hypothesis}

\begin{remark}
Hypothese~\ref{hyp:galerkin} ist numerisch bis $\varepsilon_{\mathrm{Gal}} < 10^{-12}$
verifiziert für $\lambda = 3, 5, 7, 10, 15, 20, 30, 40, 50, 75, 100$. Ein rigoroser
Beweis würde ein Galerkin-Konvergenztheorem für den spezifischen Rang-eins-Operator
\eqref{eq:rank-one-structure} erfordern, was Gegenstand laufender Arbeiten ist.
\end{remark}

\subsubsection{Massenbasierte effektive Lücke}

\begin{definition}[Effektive Cluster-Lücke]\label{def:geff}
Für $\varepsilon > 0$, sei $J_\varepsilon = \min\{J :
\sum_{j=0}^{J} |\langle q_j, k\rangle|^2 \geq 1 - \varepsilon\}$
der spektrale Massen-Schwellenwert. Die \emph{effektive Lücke} ist
\[
  g_{\mathrm{eff}}(\varepsilon) \;:=\; u_{J_\varepsilon + 1} - u_0.
\]
Der resonante Unterraum ist $V_\lambda := \operatorname{span}\{q_j : j \leq J_\varepsilon\}$.
\end{definition}

Diese Definition ersetzt die Cluster-Identifikation mit fester Toleranz
(die erratische Lücken $\sim 10^{-6}$--$10^{-4}$ produzierte) durch ein
spektrales Massen-Kriterium, das $g_{\mathrm{eff}}(\varepsilon) \approx 5$--$7$
gleichmäßig für alle getesteten $\lambda$ und alle $\varepsilon \leq 10^{-4}$ liefert.
Der zentrale Punkt: Die Quasimode-Masse $M_k(J) = \sum_{j \leq J} |\langle q_j, k\rangle|^2$
sättigt innerhalb der ersten 1--3 Eigenwerte, was eine Lücke der Ordnung eins zum Komplement erzeugt.

\subsubsection{Die drei Lemmas}

\begin{lemma}[Skalare Säkulare Aufhebung (C2ca.1)]
\label{lem:secular}
Sei $\mu = \langle \tilde{u}, (A-u_0)k\rangle$ die skalare
Komponente des Residuums entlang $\tilde{u}$. Dann
\[
  |\mu| \;\leq\; s_\lambda \;:=\; O(e^{-c \cdot N/L}).
\]
Unter Hypothese~\ref{hyp:galerkin} erzeugt die Rang-eins-Säkulargleichung
$1 + \tilde{C} \cdot m_H(w) = 0$ zwei Terme der Größenordnung $\sim 10^{-2}$, die sich zu
$\sim 3 \times 10^{-7}$ aufheben (eine Aufhebung um mehr als Faktor $200$).
\end{lemma}

\begin{proof}[Beweisskizze]
Die säkulare Identität (aus der Rang-eins-Resolventen-Formel) liefert
\[
  \frac{\mu}{\alpha} = (\bar{w} - u_0)
    + \frac{\langle f_\lambda, k_\perp\rangle}{\alpha},
\]
wobei $\alpha = \langle\tilde{u},k\rangle$, $\bar{w}$ die säkulare
Wurzel ist und $f_\lambda$ die $H$-Spektralverteilung von
$\tilde{u}$ kodiert. Beide Terme sind $O(10^{-2})$, tragen jedoch
aufgrund der Verschachtelungsstruktur $t_j \leq u_j \leq t_{j+1}$
(wobei $t_j$ die $H$-Eigenwerte sind) entgegengesetzte Vorzeichen. Die
Aufhebung wird durch den Galerkin-Trunkierungsfehler kontrolliert.
\end{proof}

\begin{lemma}[Projizierter Poisson-Quasimode (C2ca.2)]
\label{lem:projected}
Sei $h = P_0(A-u_0)k$ die $\tilde{u}$-orthogonale Komponente
des Residuums, wobei $P_0 = I - |\tilde{u}\rangle\langle\tilde{u}|$.
Dann
\[
  \|h\| \;\leq\; p_\lambda,
\]
wobei $p_\lambda \approx 2$--$3 \times 10^{-6}$ eine
Galerkin-Konstante ist (gleichmäßig in $\lambda$ bei festem $N/L$).
\end{lemma}

\begin{proof}[Beweisskizze]
Die Rand-Defekt-Zerlegung (aus der Poisson-Transfer-Theorie)
liefert $h = \frac{1}{n_{L,N}} P_0 \Pi_{L,N}(E_L^{\mathrm{bulk}} + B_L)$,
wobei $E_L^{\mathrm{bulk}}$ zu $> 99.999\%$ in $\tilde{u}$-Richtung zeigt
(sein $P_0$-Leakage ist $< 5 \times 10^{-7}$), und $B_L$ ein lokal am Rand
befindlicher Term ist, der durch den exponentiellen Abfall des
Poisson-Kerns bei $x = L$ kontrolliert wird.
\end{proof}

\begin{lemma}[Koerzives Komplement (C2ca.5)]
\label{lem:coercive}
Unter Hypothese~\ref{hyp:galerkin} existiert $g_* > 0$
(unabhängig von $\lambda$), sodass für alle
$v \in V_\lambda^\perp$ mit $\|v\| = 1$ gilt:
\[
  \langle (A^{(N)} - u_0)\, v,\, v\rangle \;\geq\; g_*.
\]
Numerisch ist $g_* \approx 5$--$7$ für alle getesteten $\lambda$.
\end{lemma}

\begin{proof}[Beweisskizze]
Durch die Rang-eins-Cauchy-Verflechtung gilt $u_j \geq t_j$ für alle $j$,
wobei $t_j$ die Eigenwerte von $H^{(N)}$ sind. Der spektrale
Massen-Schwellenwert $J_\varepsilon$ erfüllt $J_\varepsilon \leq 2$ für alle
getesteten $\lambda$, also $u_{J_\varepsilon+1} \geq t_{J_\varepsilon+1}$.
Die $H$-Spektrallücke $t_{J+1} - t_1$ ist von Ordnung eins (bestimmt durch
die kinetische Energie) und ist nach unten beschränkt durch die Lücke des
potentialfreien Laplace-Operators $\pi^2/L^2 \cdot ((J+1)^2 - 1)$, was $\geq 5$
für $J \leq 2$ und $L \leq 10$ ergibt.
\end{proof}

\subsubsection{Massenkonzentration als Korollar}

\begin{theorem}[Massenkonzentration --- Bedingter Abschluss von (MS2)]
\label{thm:mass-concentration}
Unter Hypothese~\ref{hyp:galerkin} und mit den drei obigen Lemmas:
\[
  \|(I - P_{V_\lambda})\, k_\lambda\|
  \;\leq\; \frac{r_\lambda}{g_*}
  \;\leq\; \frac{s_\lambda + p_\lambda}{g_*}
  \;\approx\; 6 \times 10^{-7},
\]
wobei $r_\lambda = s_\lambda + p_\lambda$ die gesamte Residual-Schranke
aus Lemmas~\ref{lem:secular} und~\ref{lem:projected} ist, und
$g_*$ die koerzive Lücke aus Lemma~\ref{lem:coercive} ist.
\end{theorem}

\begin{proof}
Dies ist das Standard-Argument vom Quasimode zum Projektor. Zerlege
$k = k_{\mathrm{cl}} + k_\perp$ mit $k_{\mathrm{cl}} \in V_\lambda$,
$k_\perp \in V_\lambda^\perp$. Da $V_\lambda$ $A$-invariant ist:
\begin{align*}
  g_* \|k_\perp\|^2
  &\leq \langle k_\perp, (A-u_0) k_\perp \rangle
  = \langle k_\perp, (A-u_0) k \rangle \\
  &\leq \|k_\perp\| \cdot \|(A-u_0)k\|
  = \|k_\perp\| \cdot R_0.
\end{align*}
Falls $\|k_\perp\| > 0$, dividiere durch $\|k_\perp\|$:
$\|k_\perp\| \leq R_0/g_*$. Das Residuum zerfällt als
$(A-u_0)k = \mu\tilde{u} + h$, also
$R_0 = \|(A-u_0)k\| \leq |\mu| + \|h\| \leq s_\lambda + p_\lambda
= r_\lambda$.
\end{proof}

\begin{remark}[Konvergenz gegen Null]
Bei festem $N/L = 12$ ist die Schranke $r_\lambda/g_* \approx 6 \times
10^{-7}$ bereits klein, konvergiert aber nicht gegen Null. Um
$\|(I-P_{V_\lambda})k\| \to 0$ für $\lambda \to \infty$ zu erhalten (wie
vom CCM-Programm gefordert): wähle $N = N(\lambda)$ mit $N/L \to \infty$
(z.B. $N = L^{1+\delta}$). Dann strebt $r_\lambda^{(N)} \sim
e^{-c \cdot N/L} \to 0$, während $g_* > 0$ nach unten beschränkt bleibt,
was $\|k_\perp\| \to 0$ ergibt.
\end{remark}

\begin{remark}[Numerische Verifikation]
Die Schranke in Theorem~\ref{thm:mass-concentration} ist zu 75--84\% scharf:
Die tatsächlichen Werte von $\|k_\perp\|$ liegen innerhalb eines Faktors
von $0.75$--$0.84$ der oberen Schranke $R_0/g_*$, was die Beinahe-Optimalität
der Quasimode-zu-Projektor-Abschätzung bestätigt. Vollständige Diagnosen
zur spektralen Masse (Eigenwertspektren, Überlappungskoeffizienten, kumulative
Massenkurven) sind für $\lambda = 3, \ldots, 100$ berechnet und im Begleit-Repository
verfügbar.
\end{remark}

\subsubsection{Strukturelle Zusammenfassung}

Die Beweisarchitektur lautet:
\[
\underbrace{\text{C2ca.1 (Säkular)}}_{\text{Lemma~\ref{lem:secular}}}
+ \underbrace{\text{C2ca.2 (Projiziert)}}_{\text{Lemma~\ref{lem:projected}}}
\;\longrightarrow\; R_0 \leq r_\lambda
\;\xrightarrow{\;\text{Lemma~\ref{lem:coercive} } (g_*)\;}
\;\|k_\perp\| \leq r_\lambda / g_*.
\]
Drei unabhängige Abschätzungen, eine Standardungleichung, eine bedingte
Hypothese. Die vorherige Vier-Lemma-Struktur (bei der ``Massenkonzentration''
als eigenständiges viertes Lemma aufgeführt war) ist damit hinfällig:
Das vierte Lemma ist ein Korollar, kein Input.

\subsection{RFEP Transfer-Paket für physikalische Nachträge}
\label{sec:rfep-physics-transfer}

Das Drei-Lemma-Endspiel ist nicht nur ein Instrument der Route~B. Es liefert
auch ein wiederverwendbares RFEP-Muster für die physikalische Seite des
FST-Programms. Die physikalischen Begleitpapiere sollten den Riemannschen
Beweisversuch nicht unverändert importieren; vielmehr sollten sie die Operator-Grammatik importieren:
\[
  A_X^{(N)} = H_X^{(N)} + B_X^{(N)}, \qquad
  \|(I-P_{V_X})k_X\| \leq \frac{r_X}{g_X}.
\]
Hierbei bezeichnet $X$ einen physikalischen Zweig wie Yang--Mills,
Navier--Stokes, NS-LDI, Turbulenz oder CRM; $B_X^{(N)}$ ist der
Rang-eins- oder Endlich-Rang-Defekt; $k_X$ ist der RFEP-Kandidat;
$r_X$ ist sein Residuum; und $g_X$ ist die koerzive Lücke auf dem
orthogonalen Komplement. Die Arbeitsnotiz \cite{RFEPTransferNote}
dokumentiert dies als die RFEP--Zookeeper Normalform.

\begin{proposition}[RFEP Transfer-Schema, bedingt]
\label{prop:rfep-transfer-schema}
Sei $X$ ein physikalischer FST-Nachtrag. Angenommen, dass:
\begin{enumerate}[leftmargin=*]
\item das RFEP-Funktional für $X$ einen geschlossenen quadratischen oder
  Operator-Repräsentanten $A_X$ auf einem Form-Definitionsbereich $V_X$ zulässt;
\item es eine Galerkin-Ausschöpfung $A_X^{(N)}$ gibt, die treu zum
  relevanten niedrigen Spektral-Cluster ist;
\item $A_X^{(N)}$ zerfällt in einen koerziven Hintergrund $H_X^{(N)}$ plus einen
  Rang-eins- oder Endlich-Rang-Defekt $B_X^{(N)}$;
\item der RFEP-Kandidat $k_X$ ein Residuum
  $\|(A_X^{(N)}-u_0)k_X\|\le r_X=s_X+p_X$ aufweist, das einen skalaren säkularen
  Teil und einen projizierten Quasimode-Teil besitzt;
\item auf $V_X^\perp$ eine koerzive Komplement-Lücke $g_X>0$ existiert.
\end{enumerate}
Dann gilt
\[
  \|(I-P_{V_X})k_X\|\le r_X/g_X.
\]
Wenn der von $P_{V_X}$ ausgewählte Cluster einfach oder RFEP-selektiv ist,
ist der stabile physikalische Zustand innerhalb des entsprechenden
Galerkin-Sektors eindeutig.
\end{proposition}

\begin{proof}
Der Beweis ist identisch zu Theorem~\ref{thm:mass-concentration}:
Zerlege $k_X$ in Cluster- und Komplement-Komponenten, verwende Koerzivität
auf dem Komplement und beschränke die Komplement-Komponente durch das Residuum.
Es folgt keine physik-spezifische Schlussfolgerung, solange die fünf
Hypothesen nicht im gegebenen Bereich verifiziert sind.
\end{proof}

\begin{remark}[Was RFEP gewinnt]
\label{rem:rfep-gain}
RFEP gewinnt eine konkrete Beweis-Checkliste. Ein Physik-Paper kann nun fragen:
Was ist der Operator oder die Form $A_X$; was ist der Defekt $B_X$; was
besagt die Aussage zur Galerkin-Treue; wo ist die Komplement-Lücke;
und wie bewahrt der Limes des Definitionsbereichs die RFEP-Selektion? Das ist
stärker als eine erzählerische Analogie, aber schwächer als ein Theorem in der jeweiligen Domäne.
\end{remark}

\begin{remark}[Lesarten in den Domänen]
\label{rem:domain-readings}
Für Yang--Mills ist die angestrebte Lesart die Massenlücke oder Gibbs-Eindeutigkeit
als koerzives Komplement. Für Navier--Stokes und NS-LDI sind es dissipative Selektion
und Log-Distanz-Defekt-Kontrolle als Aussage über einen stabilen Cluster. Für Turbulenz
ist es das K41-Spektrum als massenkonzentrierter RFEP-Minimierer. Für CRM bleibt
die Hauptbeweissprache strikte Konvexität/UCU, aber das Transfer-Paket liefert die
Operatorsprache für perturbative oder randgekoppelte Varianten. Das Prime-Hub-Programm
\cite{PrimeHubSynthesis} ist die natürliche intermediäre Brücke: Frontier-Prime-Einfügungen
und Bouquet-Trunkierungen kommen Endlich-Rang-Rand-Defekten bereits sehr nahe.
\end{remark}

\subsection{Die Sonin-Raum-Brücke (Kontext aus CCM 2024)}
\label{sec:sonin}

Die Operator-Norm-Schranke \eqref{eq:dynamic-isometry} ist eingebettet in den
analytischen Kontext der semilokalen Connes-Consani-Moscovici-Spurformel
\cite{ConnesConsaniMoscovici2024}:

\begin{theorem}[CCM 2024, semilokale Zerlegung]\label{thm:ccm-semilocal}
Im archimedischen Fall lässt der Prolat-Operator die Zerlegung
\[
  PW_\lambda \;=\; D^{(\lambda)}_{\log}{}^2 \;+\; \Gamma,
\]
zu, wobei $\Gamma$ ein Graduierungs-Operator auf orthogonalen Polynomen ist.
Dementsprechend zerfällt sein Spektrum in
\begin{itemize}[leftmargin=*]
\item einen \emph{positiven Teil}, der tiefliegende Nullstellen von
  $\zeta(1/2+is)$ realisiert (Infrarot- / Weil-Regime),
\item einen \emph{negativen Teil} auf dem Sonin-Raum, der die UV-Asymptotik
  der Nullstellenverteilung realisiert (Ultraviolett- / Connes-Moscovici 2022-Regime).
\end{itemize}
\end{theorem}

Unter dieser Zerlegung stimmt die Weilsche Form $QW_\lambda$ mit der Einschränkung
von $PW_\lambda$ auf den positiven spektralen Unterraum überein, bis auf die
Rang-eins-Korrektur, die durch den Spurformel-Fehler $W_\infty - W_0$ kodiert wird.
Dies identifiziert die Operatordifferenz in \eqref{eq:dynamic-isometry} mit dem
arithmetisch-archimedischen Defekt nach Connes-Consani \cite{ConnesConsaniWeilPositivity2020}.
Ob dieser Defekt die Art von Operator-Norm-Schranke zulässt, die in
Entwurf~\ref{bp:fst-transfer} gefordert wird, ist eine offene Frage.

% =====================================================
\section{Die Wick-Mellin-Rotation und ihre drei Hindernisse (Route A)}
\label{sec:wick-mellin-obstacles}
% =====================================================

\subsection{Die Rotation ist keine gewöhnliche Wick-Rotation}

Die Standard-Wick-Rotation verbindet $\Re\beta > 0$ (thermisch) mit
$\Im\beta \in \R$ (unitär). Was Route~A benötigt, ist anders:
$\Re\beta > 1$ (BC-Gibbs) zu $\Re\beta = 1/2$ (Meyers kritische Linie),
d.h. eine \emph{Halbebenen-Fortsetzung durch den Mangoldt-Pol bei
$\beta = 1$}.

Der Gibbs-Operator $e^{-\beta H_{BC}}$ ist:
\begin{itemize}
\item Spurklasse für $\Re\beta > 1$.
\item Beschränkt, aber nicht Spurklasse für $0 < \Re\beta \leq 1$.
\end{itemize}

Auf der kritischen Linie existiert der klassische Gibbs-Zustand nicht;
nur die meromorph fortgesetzte Funktion $\zeta(\beta)$ bleibt bestehen.
Drei konkrete Hindernisse strukturieren die offene technische Arbeit.

\subsection{Hindernis (O1): Die Distributionsklasse}

Gesucht ist ein präziser Fréchet-Raum $\V_\chifn$, dessen Dualraum
Meyers Eigendistributionen enthält.

\begin{definition}[der Meyer-Raum]
\label{def:meyer-space}
$\V_\chifn := \S_0(\A^\times)^\chifn / \K_{\Q^\times}^\chifn$,
wobei $\S_0 = \{f : Z(f, s)$ ganz$\}$ ist und $\K_{\Q^\times}$ den
Abschluss der $\Q^\times$-Koinvarianten bezeichnet.
\end{definition}

\textbf{Halbnormen:} archimedisch gewichtete Schwartz-Normen
$\|f_\infty\|_{n,m} = \sup |x^m \partial^n f_\infty(x)|$; p-adische
Normen $\|f_p\|_N^p$ auf kompakten Teilmengen von $\Q_p^\times$.
$\V_\chifn$ ist ein Fréchet-Raum in der Quotiententopologie.

\textbf{$\Q^\times$-Invarianz der Mellin-Auswertung.} Für $f \in
\S_0$, $q \in \Q^\times$:
\[
  v_\gamma^{(pre)}(T_q f) = |q|^{1/2+i\gamma} \cdot
    v_\gamma^{(pre)}(f) = v_\gamma^{(pre)}(f)
\]
nach der Produktformel $|q| = 1$. Genau hier tritt (FST-7A)(ii) ein.

\textbf{Status von (O1).} Die Distributionsklasse ist eng mit
anerkannter adelischer Analysis und Meyer 2005 verbunden; das Bild aus
Fréchet-Quotient und Mellin-Auswertung ist mathematisch wiedererkennbar.
Der eigentliche analytische Gehalt von (O1) ist die Dichte/Vollständigkeit
der Mellin-Auswertungen im Dualraum von $\V_\chifn$ --- im Kern also ein
Vollständigkeitssatz, der aus Meyer 2005, Abschnitte~4--5, importiert oder
in diesem Stil bewiesen werden muss. Relativ zu (O2) und (O3) ist dies
vergleichsweise solide, aber die abstrakte Verallgemeinerung bleibt
nichttrivial.

\subsection{Hindernis (O2): Die Fourier-Dualität}

Gesucht ist eine explizite operatoralgebraische Fourier-Transformation
\[
  \mathcal F : \HH_\Q^{\mathrm{Sch}} \to \V_{\chifn_0}
\]
zwischen der Bost-Connes-Hecke-Algebra (Schwartz-Unteralgebra) und
Meyers Fréchet-Raum.

\textbf{Vier Brückenaussagen (Ziele, keine Sätze):}

\begin{itemize}
\item (F1) Elementkorrespondenz: $\mu_n \leftrightarrow
  \mathbf{1}_n$.
\item (F2) Zustandskorrespondenz: KMS-Zustand $\phi_{\beta,\sigma}
  \leftrightarrow$ Galois-verdrehtes Tate-Funktional.
\item (F3) Zeitentwicklung: $\alpha_t^{BC}$ $\leftrightarrow$
  Mellin-Verschiebung durch analytische Fortsetzung
  $t = i(\beta - 1/2)$.
\item (F4) Spektrenabgleich: $\{\log n\}$ $\leftrightarrow$
  Kontinuumspektrum mit singulären Punkten bei $\{\gamma_\rho\}$.
\end{itemize}

\textbf{Kandidatenkonstruktion über Galois-Mittelung:}
\[
  \mathcal F(a)([f]) := \int_{\hat\Z^\times}
    \phi_{\infty, \sigma}(a) \cdot \sigma(f) \, d\sigma.
\]

\textbf{Status von (O2).} Eine konkrete Zielabbildung ist beschrieben,
aber die notizartigen Aussagen (F1)--(F4) enthalten noch keine Beweise
für Existenz, Stetigkeit, Injektivität und Surjektivität. Die Grenz-KMS-
Zustände und ihre Stetigkeit sind nicht konstruiert. Die vorgeschlagene
Fourier-Abbildung bleibt genau an der Stelle heuristisch, an der die
Fortsetzung auf die kritische Linie eintreten müsste. (O2) sollte daher
als eigenständiges Forschungsproblem in der Operatoralgebra behandelt
werden, nicht als schnelles Brückenlemma.

\subsection{Hindernis (O3): Konvergenzkontrolle}

Der Mangoldt-Pol bei $\beta = 1$ ist ein Hindernis für direkte
Operatorfortsetzung. Ein Renormalisierungsansatz liegt nahe, aber --- wie
das Audit zutreffend betont --- sind weder die Spurklasse-Eigenschaften
noch die analytische Fortsetzung des renormalisierten Operators derzeit
durch strenge Argumente gerechtfertigt.

\begin{definition}[renormalisierter Gibbs-Operator, vorläufig]
\label{def:ren-gibbs}
\[
  G_\beta^{\mathrm{ren}} := e^{-\beta H_{BC}}
    - \frac{\Pi_0}{\beta - 1},
\]
wobei $\Pi_0$ der kanonische Projektor auf den
``$n = 1$''-Identitätssektor der Hecke-Algebra ist (der den
Mangoldt-Pol trägt).
\end{definition}

\begin{conjecture}[Spurklasse-Fortsetzung durch den Pol,
\textbf{offen})]
\label{conj:ren-trace-class}
$G_\beta^{\mathrm{ren}}$ ist für alle $\Re\beta > 0$ Spurklasse
(einschließlich der kritischen Linie) mit
$\Tr(G_\beta^{\mathrm{ren}}) = Z_{\mathrm{ren}}(\beta) = \zeta(\beta)
- (\beta - 1)^{-1}$.
\end{conjecture}

\begin{remark}[Rücknahme der Proposition aus v0.6]
\label{rem:o3-retraction}
In v0.6 wurde Vermutung~\ref{conj:ren-trace-class} als Proposition
formuliert. Das war falsch: Das Subtrahieren eines Pols auf Ebene der
Spuren erzeugt nicht automatisch eine Operatorfamilie der Spurklasse.
Ein operatortheoretischer Beweis erfordert mindestens eine explizite
Analyse der Rang-eins-Subtraktion auf dem Hilbertraum der BC-Darstellung,
Kontrolle der Operatornorm-Grenzwerte für $\Re\beta \searrow 1$ und eine
strenge Identifikation der fortgesetzten Familie auf $\Re\beta \le 1$.
Nichts davon folgt aus der skalaren Zetaregularisierung. Wir nehmen die
Behauptung aus v0.6 zurück und markieren Vermutung~\ref{conj:ren-trace-class}
ausdrücklich als offen.
\end{remark}

\textbf{FST-Interpretation.} Die Subtraktion des Mangoldt-Pols erinnert
an eine RG-Operation, die Bulk-Freie-Energie vom RFEP-stabilen Anteil
trennt. Die verbleibende Polstruktur $\{\beta = \rho\}$ würde, vermutungsweise,
Meyers Eigendistributionen als RG-kritische Modi eines renormalisierten
Systems liefern. Dies bleibt ein strukturelles Bild, kein Satz.

\subsection{Status der drei Hindernisse}

Wir behaupten \emph{nicht}, dass (O1)--(O3) auf Standardtechniken
reduzierbar sind. Die ehrliche Einschätzung lautet:
\begin{itemize}[leftmargin=*]
\item (O1) ist vergleichsweise solide --- stilistisch nahe an Meyer 2005
  und klassischer adelischer Analysis; die abstrakte Verallgemeinerung ist
  nichttrivial, aber plausibel.
\item (O2) ist offen, aber strukturiert: Eine konkrete Zielabbildung ist
  angegeben, doch die analytischen Bausteine (Stetigkeit, Grenz-KMS-
  Zustände, Injektivität, Surjektivität) fehlen.
\item (O3) ist in seiner aktuellen Form tief und stark spekulativ. Die
  Spurklasse-Fortsetzung durch den Pol wurde in früheren Versionen zu
  stark behauptet; siehe Bemerkung~\ref{rem:o3-retraction}.
\end{itemize}
Der Satz ``alle drei Hindernisse reduzieren sich auf Standardtechniken'',
der in v0.6 stand, ist nicht korrekt und wird zurückgenommen.

% =====================================================
\section{Nichtabelsche Erweiterung: Die Selberg-Kompatibilitätsprüfung}
\label{sec:selberg}
% =====================================================

\subsection{Warum eine Kompatibilitätsprüfung wichtig ist}

Die FST-Basisaxiome (Abschnitt~\ref{sec:axioms}) verlangen, dass $G$
abelsch ist. Selbergs Beweis der RH für $Z_\Gamma$ ist ein
\emph{klassischer} nichtabelscher Fall. Wenn eine natürliche nichtabelsche
Erweiterung der FST-Axiome mit Selbergs Setting \emph{kompatibel} ist ---
also wenn sich die Axiome widerspruchsfrei in das nichtabelsche Setting
übersetzen lassen und die Selberg-Strukturen in die übersetzten Axiome
passen ---, dann ist dies Evidenz dafür, dass die Axiome keinen offensichtlich
nur abelschen strukturellen Bias haben.

Wir betonen: Eine Kompatibilitätsprüfung ist kein Beweis. Der nichtabelsche
FST-Fixpunktsatz (Vermutung~\ref{conj:nonabelian-fixed-point} unten) ist
eine Vermutung, kein Satz, der Selberg aus den neuen Axiomen zurückgewinnt.

\subsection{Selberg-Rekapitulation}

Für $\Gamma \subset \mathrm{PSL}_2(\R)$ kokompakt und
$\mathcal F = \Gamma \backslash \mathbb H$:
\begin{itemize}
\item $\Delta_{\mathcal F}$ ist selbstadjungiert auf $L^2(\mathcal F)$,
  mit diskretem Spektrum $\{\lambda_n\}$.
\item Selberg-Zeta $Z_\Gamma(s)$: nichttriviale Nullstellen
  $\rho_n = 1/2 + i r_n$ mit $\lambda_n = 1/4 + r_n^2$.
\item Selberg-RH: $r_n \in \R$ folgt automatisch aus
  $\lambda_n \geq 0$ (Selbstadjungiertheit des Laplace-Operators).
\end{itemize}

Die Selberg-Spurformel verbindet $\sum_n h(r_n)$ mit Daten geschlossener
Geodäten --- eine nichtabelsche Formel vom Poisson-Weil-Typ.

\subsection{Axiomprüfung für Selberg}

\begin{table}[h]
\centering
\small
\begin{tabular}{@{}lcl@{}}
\toprule
\textbf{Axiom} & \textbf{Status} & \textbf{Kommentar} \\
\midrule
(G) abelsch & nein & verlangt nichtabelsche Erweiterung \\
(R) Renormalisierung & $\checkmark$ & geodätischer Fluss \\
(N) arithmetischer Nukleus & partiell & tannakisches Analogon \\
($\chifn$) Charakter & $\checkmark$ & ersetzt durch Darstellung \\
(P1--P3) Muster A & partiell & spektrale Zerlegung existiert \\
(SGE) Gruppenstruktur & $\checkmark$ & $\Gamma$ ist eine Gruppe \\
(DS1--DS3) & $\checkmark$ & geodätische Dynamik \\
(RFEP) & $\checkmark$ & Laplace-Funktional \\
\bottomrule
\end{tabular}
\caption{FST-Axiomprüfung für das Selberg-Setting. Das abelsche
Axiom (G) scheitert; alle anderen sind mit einer natürlichen
nichtabelschen Erweiterung \emph{kompatibel}.}
\label{tab:selberg-axioms}
\end{table}

\subsection{Tannakische Erweiterung (Skizze)}

Ersetze:
\begin{itemize}
\item die abelsche Gruppe $G$ durch die tannakische Kategorie $\T$ der
  Darstellungen.
\item den Charakter $\chifn$ durch eine irreduzible Darstellung
  $\pi \in \T$.
\item Pontryagin-Dualität durch tannakische Rekonstruktion.
\end{itemize}

Für Selberg entsteht $\T$ aus $\Rep(\mathrm{PSL}_2(\R))$, eingeschränkt
auf $\Gamma$-invariante Untermoduln.

\begin{conjecture}[nichtabelscher FST-Fixpunktsatz]
\label{conj:nonabelian-fixed-point}
Unter einer geeigneten nichtabelschen (tannakischen) Formulierung der
FST-Axiome und einem nichtabelschen Analogon (FST-7') existiert ein
Operator $H_\Gamma$ auf einem Fréchet-Raum $\V_\Gamma$, dessen
Punktspektrum den Selberg-Nullstellenordinaten entspricht und der für die
Petersson-Paarung $B$-symmetrisch ist. Die Realität des Spektrums würde aus
Selbstadjungiertheit folgen (durch eine automatische Laplace-Struktur),
mithin RH für $Z_\Gamma$.
\end{conjecture}

\subsection{Was die Selberg-Kompatibilitätsprüfung etabliert}

\begin{enumerate}[leftmargin=*]
\item \textbf{Strukturelle Kompatibilität.} FST (in nichtabelscher Form)
  ist nicht offensichtlich inkompatibel mit einem bekannten RH-Beweis ---
  ein Plausibilitätscheck der Axiome.
\item \textbf{Gestalt einer nichtabelschen Erweiterung.} Das tannakische
  Vokabular liefert eine natürliche Sprache für nichtabelsche
  darstellungstheoretische Inhalte.
\end{enumerate}

Wir behaupten \emph{nicht}: dass die nichtabelschen FST-Axiome Selbergs
Beweis \emph{neu herleiten}, dass (FST-7'(C)) aus Anosov-Ergodizität in
einer Weise folgt, die als Satz des tannakischen FST-Rahmens formalisiert
ist, oder dass das nichtabelsche Axiompaket eine kategorientheoretische
Reife besitzt, die mit dem abelschen Paket vergleichbar wäre. Gezeigt wird
eine Analogie und ein Kompatibilitätstest, kein bewiesener nichtabelscher
Fixpunktsatz.

\subsection{Die Riemann-vs.-Selberg-Diagnose}

Der entscheidende Unterschied:
\begin{itemize}
\item Selberg: Das arithmetisch-adelische Analogon von (FST-7) ist
  \emph{geometrisch natürlich} --- der Laplace-Operator entsteht aus
  riemannscher metrischer Geometrie.
\item Riemann: (FST-7) ist \emph{arithmetisch adelisch} --- ohne analoge
  Geometrie; Meyers Fréchet-distributionale Struktur spielt die analoge
  strukturelle Rolle.
\end{itemize}

Der thermodynamische Bost-Connes-Mechanismus (Abschnitt~\ref{sec:emergence})
wurde in früheren Entwürfen als dynamische Brücke über diese Lücke
vermutet; wir betrachten dies nun als offen.

% =====================================================
\section{Klassifikation über FST-Struktur}
\label{sec:classification}
% =====================================================

\begin{table}[h]
\centering
\small
\begin{tabular}{@{}llcccl@{}}
\toprule
\textbf{Familie} & \textbf{Typ} & \textbf{SGE} & \textbf{(FST-7)} & \textbf{FST-Mechanismus} & \textbf{Status} \\
\midrule
Riemann $\zeta(s)$ & abelsch & $\checkmark$ & $\checkmark$ (Tate) & Meyer / CCM-Prolat & offen; CCM-Numerik stark \\
Dirichlet $L(s, \chifn)$ & abelsch & $\checkmark$ & $\checkmark$ (Tate) & Meyer-adelisch & offen \\
Hecke $L(s, \chifn_K)$ & abelsch & $\checkmark$ & $\checkmark$ (Tate-Weil) & Meyer-körperadelisch & offen \\
Selberg $Z_\Gamma$ & nichtabelsch & $\checkmark$ & $\checkmark$ (Spurformel) & Laplace-direkt & bewiesen (Selberg) \\
Ihara-Petersen & endlich & $\checkmark$ & N/A (endlich) & trivial & bewiesen (Bass-Hashimoto) \\
Artin (nichtabelsch) & nichtabelsch & $\checkmark$ & offen & Tannaka-Langlands & offen \\
Automorphe $L$ auf $GL_n$ & nichtabelsch & $\checkmark$ & offen & Langlands-funktoriell & offen \\
$\F_1$-hypothetisch & monoidal & ? & offen & Connes-Consani & offen \\
Prime-Hub (OP5) & unklar & ? & offen & keine natürliche Gruppe & offen \\
\bottomrule
\end{tabular}
\caption{Klassifikation zetatypischer Familien unter FST-Axiomen,
mit einer tannakischen Erweiterung für nichtabelsche Fälle. FST fungiert
hier als Ordnungsrahmen, nicht als Beweismaschine.}
\label{tab:classification-v2}
\end{table}

\textbf{Beobachtung.} FST liefert eine prognoseartig wirkende Klassifikation:
Familien, in denen (FST-7) gilt, erlauben Meyer-artige Konstruktionen; die
nichtabelschen Fälle verlangen eine tannakische Erweiterung (SGE liefert
weiterhin die Gruppenstruktur). Prime-Hub ist ungelöst, weil bislang keine
natürliche Gruppe identifiziert wurde.

% =====================================================
\section{Das Emergenz-Programm (Route C): Diagnose eines tiefen Hindernisses}
\label{sec:emergence}
% =====================================================

Frühe Entwürfe dieses Projekts schlugen Route~C vor: einen thermodynamischen
Mechanismus, durch den FST zusammen mit zwei zusätzlichen Selektionsaxiomen
und Cornelissen-Marcolli-Starrheit (FST-7) als Satz ``ableiten'' sollte.
Die spätere Analyse (vgl.~\cite{L1KAttack,EmergenceAttack}) zeigt, dass
diese Route die arithmetische Tiefe eher \emph{verlagert} als auflöst.
Dieser Abschnitt stellt das Emergenz-Programm daher in seiner nun
diagnostischen Form dar: als Identifikation des Ortes, an dem der schwere
arithmetische Input liegt, nicht als fast vollständigen Pfad zu (FST-7).

\subsection{Zwei zusätzliche FST-Axiome als \emph{Selektionsprinzipien}}
\label{sec:extra-axioms}

\begin{definition}[FST-Free Axiom]\label{def:fst-free}
Die multiplikative Halbgruppe $N \le G$ ist eine freie kommutative
Halbgruppe, erzeugt von irreduziblen ``Primzahlen'' $\{p_i\}_{i \in I}$.
\end{definition}

\begin{definition}[FST-Min Axiom]\label{def:fst-min}
Unter allen FST-Strukturen, die die Basis-Axiome plus
(FST-Free) erfüllen und einen kritischen Phasenübergang zulassen,
wählt ein zusätzliches Auswahlkriterium jene Struktur aus, deren zugrunde liegender
globaler Körper $K$ das \emph{Anfangsobjekt} in der Kategorie der globalen
Körper der Charakteristik null ist. Im abelschen Charakteristik-null-Fall
wählt dies $K = \Q$ aus.
\end{definition}

\begin{remark}[Ehrliche Lesart]
\label{rem:fst-free-min-imports}
(FST-Free) ist eine multiplikative Freiheits-Annahme an $N$. (FST-Min)
ist ein Selektionskriterium minimaler Beschreibung. Beide sind \emph{starke
Selektionsprinzipien, die arithmetischen Inhalt kodieren} --- sie sind
keine Konsequenzen der Basis-Axiome. Insbesondere fungieren sie in der
untenstehenden Reduktion als externe arithmetische Importe, nicht
als abgeleitete Fakten.
\end{remark}

\subsection{Die Cornelissen-Marcolli Brücke}
\label{sec:cm-bridge}

\begin{theorem}[Cornelissen-Marcolli
2010 \cite{CornelissenMarcolli2010}]\label{thm:cm-rigidity}
Seien $K_1, K_2$ Zahlkörper und seien $(\mathcal A_{K_i},
\sigma_{K_i})$ ihre zugehörigen Bost-Connes-Systeme. Die folgenden Aussagen
sind äquivalent:
\begin{enumerate}[label=(\roman*)]
\item $K_1 \cong K_2$ als Körper.
\item $(\mathcal A_{K_1}, \sigma_{K_1}) \cong (\mathcal A_{K_2},
  \sigma_{K_2})$ als $C^\ast$-dynamische Systeme.
\end{enumerate}
Insbesondere determiniert das BC-System den zugrunde liegenden Zahlkörper
bis auf Isomorphie.
\end{theorem}

\begin{remark}
Theorem~\ref{thm:cm-rigidity} behauptet \emph{nicht}, dass die
Zustandssumme (Partitionsfunktion) allein den Körper determiniert (zwei
nicht-isomorphe Körper können dieselbe Dedekindsche Zetafunktion teilen).
Die volle $C^\ast$-Algebra plus die einparametrige Gruppe wird benötigt.
Dies ist eine wichtige Hypothese des Theorems; sie ist stärker als der Input, den
das Emergenz-Programm tatsächlich liefert.
\end{remark}

\subsection{Thermodynamische Bedingungen}
\label{sec:thermo-B}

Folgend \cite{LemmaThreeTwo,ThermoDerivation} extrahiert man aus
FST$_0$ + RFEP + (FST-Free) vier \emph{thermodynamische Bedingungen} an
das zugehörige QSM-System $(\mathcal A, \sigma)$:

\begin{description}[leftmargin=*,style=nextline]
\item[(B1)] \textbf{Dirichlet-Struktur.} $Z(\beta) = \sum_{n \in X}
  w_n\, n^{-\beta}$ für eine abzählbare Indexmenge $X$ und positive
  Gewichte $w_n$.
\item[(B2)] \textbf{Euler-Produkt.} $Z(\beta) = \prod_{p \in P}
  (1-w_p p^{-\beta})^{-1}$ für eine atomare Erzeugermenge $P$
  (``Primzahlen'').
\item[(B3)] \textbf{Einfacher Pol bei $\beta=1$.} $\lim_{\beta \to 1^+}
  (\beta-1)\, Z(\beta) = c > 0$.
\item[(B4)] \textbf{FST$_0$-Konformität.} $\mathcal A$ ist das
  universelle Muster-A verschränkte Produkt (crossed product) von $N$
  mit seiner einparametrigen Skalierung.
\end{description}

\begin{proposition}[Thermodynamische Ableitung, partiell]\label{prop:thermo-derivation}
Unter den FST$_0$-Basis-Axiomen zusammen mit (FST-Free) und RFEP mit
einer einparametrigen thermodynamischen Variable gibt es eine heuristische
Route zu (B1), (B2) und (B4); (B3) erfordert einen weiteren arithmetischen Input
(Identifikation des thermodynamischen Pols mit dem Dedekind-Pol).
\end{proposition}

\begin{remark}[Status von Proposition~\ref{prop:thermo-derivation}]
\label{rem:thermo-status}
Die Beweisskizzen in v0.4--v0.6 präsentierten dies als eine geradlinige
Ableitung. Die ehrliche Position ist: (B1), (B2), (B4) folgen aus
RFEP + (FST-Free) durch Standard-Argumente vom Typ kritischer Exponenten;
(B3) --- welches der inhaltliche arithmetische Punkt ist --- ist der Ort, an dem die
Identifikation mit einem echten Zahlkörper-Dedekind-Pol ins Spiel kommt.
Dies folgt nicht allein aus der thermodynamischen Seite; es erfordert den Input,
der in der Lückenanalyse von \cite{L1KAttack} entwickelt wurde.
\end{remark}

\subsection{Das ehemalige ``Hauptreduktions-Theorem''}
\label{sec:former-reduction}

Frühere Versionen dieses Papers formulierten ein Theorem (``Reduktion der
Emergenz-Hypothese''), das behauptete, FST$_0$ + (FST-Free) +
(FST-Min) + RFEP + Cornelissen-Marcolli-Starrheit erzwingen, dass das
zugrundeliegende QSM-System das BC-System von $\Q$ ist, und bedingen somit
(FST-7)(A) via Tate. Die aktualisierte Position ist, dass diese Kette einen
entscheidenden offenen Schritt bei (B3) hat, und grundlegender, dass die
Identifikation eines QSM-Systems $(\mathcal A, \sigma)$, das
(B1)--(B4) erfüllt, mit dem BC-System eines \emph{Zahlkörpers} $K$ \emph{mehr}
erfordert als ein Dirichlet/Euler/Pol-Muster: Es erfordert die Rekonstruktion
der additiven Struktur von $K$ aus Daten, die prima facie nur multiplikativ sind.

Wir formulieren dies formell:

\begin{conjecture}[L1-K: Rekonstruktion der additiven Struktur]
\label{conj:l1k}
Ein QSM-System $(\mathcal A, \sigma)$, das (B1)--(B4) plus
(FST-Min) erfüllt, ist isomorph als ein $C^\ast$-dynamisches System zum
Bost-Connes-System eines Zahlkörpers $K = \Q$.
\end{conjecture}

\begin{remark}[Warum L1-K schwer ist]\label{rem:l1k-hard}
Cornelissen-Marcolli charakterisieren einen Zahlkörper durch sein BC-System,
aber die Hypothesen dieses Theorems beinhalten die volle
$C^\ast$-dynamische Struktur eines BC-Systems --- insbesondere jene Teile,
die die Ringstruktur des Ganzheitsrings kodieren, einschließlich der Addition.
Die thermodynamischen Bedingungen (B1)--(B4) sind primär multiplikativ.
Diese Lücke zu schließen, ist klassischerweise äquivalent dazu, eine additive
Struktur aufzuzwingen, die kompatibel mit der multiplikativen ist, d.h. das Problem
der additiv-multiplikativen Kopplung, das vom Connes-Consani Arithmetic-Site /
$\F_1$-Programm adressiert wird. Keine Version des in diesem Projekt entwickelten
Emergenz-Programms schließt diese Lücke. Innerhalb unserer formalen Sprache
entspräche dies der Einführung noch eines weiteren Axioms, \emph{(FST-Add)},
welches die additive Struktur von $K$ kodiert --- in welchem Fall das Argument
offensichtlich zirkulär wird: wir setzen das voraus, was wir ableiten wollen.
\end{remark}

\subsection{Entwurf für Route C, in diagnostischer Form}

\begin{blueprint}[Route C, in diagnostischer Form]
\label{bp:route-c}
Unter FST$_0$ + (FST-Free) + (FST-Min) + RFEP + (B3), \emph{und
bedingt durch Vermutung~\ref{conj:l1k} (L1-K)}, schließt sich die Kette
\begin{center}
FST Daten $\;\to\;$ QSM-System, das (B1)--(B4) erfüllt
$\;\xrightarrow{\text{CM Starrheit}}\;$ BC-System von $\Q$
$\;\xrightarrow{\text{Tate}}\;$ (FST-7)(A)
\end{center}
In diesem Szenario wäre (FST-7) ein Theorem. In der Praxis ist
Vermutung~\ref{conj:l1k} das additiv-multiplikative Kopplungsproblem
des $\F_1$-Programms; das Emergenz-Programm schließt die Schleife daher
\emph{nicht} autonom.
\end{blueprint}

\subsection{Herabstufung von Route C}

\begin{remark}[Herabstufung]
\label{rem:route-c-demotion}
Folglich stufen wir Route~C von einem potenziellen ``Hauptprogramm'' zu einer
diagnostischen Nebenuntersuchung herab. Ihr Beitrag liegt in der strukturellen Klärung:
Indem man der Emergenz-Kette folgt, kann man die additiv-multiplikative
Kopplung als das einzige tiefe Hindernis identifizieren, und man sieht explizit, dass
BC-Thermodynamik allein keine additive Struktur liefert. Dies ist nützlich als ein
diagnostisches Ergebnis darüber, wo die arithmetische Tiefe liegt; es ist
als Beweiswerkzeug \emph{nicht} nützlich.
\end{remark}

\subsection{Die drei parallelen Routen unter dem FST-Schirm}
\label{sec:three-routes}

Mit der Herabstufung von Route~C stellt sich das Routenbild wie folgt dar:

\begin{center}
\small
\begin{tabular}{@{}lllll@{}}
\toprule
\textbf{Route} & \textbf{Mechanismus} & \textbf{(FST-7) Sichtweise} &
\textbf{Schlüsselobjekt} & \textbf{Status} \\
\midrule
(A) Meyer-Tate      & adelische Selbstdualität & Pontryagin        &
$\A_\Q$, $\Q$-Gitter & Kontext; (O1)--(O3) offen \\
(B) CCM-Prolat      & Zeit-Frequenz-Optimalität & Nyquist-Shannon &
$D^{(\lambda,N)}_{\log}$, $PW_\lambda$ & \textbf{Hauptprogramm; 4 Meilensteine} \\
(C) BC-Emergenz     & thermodynamischer Phasenübergang & arithm. ergodisch &
BC-Algebra            & herabgestuft (L1-K = add/mult-Kopplung) \\
\bottomrule
\end{tabular}
\end{center}

Nach Theorem~\ref{thm:fst7-equiv} sind die drei Sichtweisen von (FST-7)
als strukturelle Bedingungen äquivalent, so dass jede Route, wenn sie abgeschlossen
würde, (FST-7) und damit (in Kombination mit der Meyer-Konstruktion
oder ihrem Analogon) die Riemannsche Vermutung nach sich ziehen würde.
Der Punkt ist, dass ein Abschluss auf den drei Routen sehr unterschiedliche Dinge bedeutet, und
derzeit nur Route~B als ein aktives technisches Programm verfolgt wird.

\begin{remark}[Rolle der FST in v1.8]\label{rem:fst-role-v07}
FST konkurriert nicht mit irgendeiner der drei Routen. Sie bietet einen
gemeinsamen \emph{Klassifizierungs- und Vergleichsrahmen}: Dasselbe
(FST-7)-Axiom wird in drei äquivalenten strukturellen Formen durch drei
verschiedene technische Programme angegriffen. Der tiefe Inhalt ist die
strukturelle \emph{Äquivalenz} --- dass adelische Selbstdualität,
Zeit-Frequenz-Optimalität und arithmetische Ergodizität drei
strukturelle Gesichter eines einzigen Phänomens sind. FST beweist
für sich allein nichts über RH.
\end{remark}

% =====================================================
\section{Offene Forschungsfragen}
\label{sec:ofq}
% =====================================================

\begin{description}[style=nextline,leftmargin=*,labelwidth=3.5em]
\item[(OFQ1)] \textbf{(FST-7) aus RFEP.} Offen; die Emergenz-Route
  identifiziert (L1-K), die additiv-multiplikative Kopplung, als das
  Kernhindernis (Abschnitt~\ref{sec:emergence}).
\item[(OFQ2)] \textbf{Auflösung von (O1)--(O3).} Technische
  Vollendung der Wick-Mellin-Rotation; (O3) ist das schwerste Stück,
  und sein Spurklasse-Status ist offen (Vermutung~\ref{conj:ren-trace-class}).
\item[(OFQ3)] \textbf{Nicht-abelsche FST.} Tannakische Axiome als ein
  Vokabular für Artin- und automorphe $L$-Funktionen;
  Vermutung~\ref{conj:nonabelian-fixed-point} offen;
  vgl. \cite{TannakianAxioms} für explorative Formalisierung.
\item[(OFQ4)] \textbf{Selberg-Kompatibilität.} Kompatibilität
  etabliert (Abschnitt~\ref{sec:selberg}); ein vollständiger nicht-abelscher
  Fixpunktsatz bleibt offen.
\item[(OFQ5)] \textbf{$\F_1$-FST und die additiv-multiplikative Kopplung.}
  Das (L1-K)-Hindernis aus Abschnitt~\ref{sec:emergence} fällt mit dem
  klassischen Connes-Consani-Arithmetic-Site-Programm zusammen;
  dies ist der langfristige Rahmen, in dem Route~C möglicherweise
  wiederbelebt werden könnte.
\item[(OFQ6)] \textbf{Physikalische Realisierung.} Bost-Connes in der
  kondensierten Materie; experimenteller Phasenübergang (spekulativ).
\item[(OFQ7)] \textbf{Starrheit (Rigidity).} Adelische Struktur:
  Margulis-artige Starrheit oder Moduli? (explorativ).
\end{description}

% =====================================================
\section{Status, Ehrliche Einschätzung und Ausblick}
\label{sec:status}
% =====================================================

\subsection{Was dieses Paper etabliert}

\begin{enumerate}[leftmargin=*]
\item Axiomatisches FST-Framework zur Spektrum-Nullstellen-Identifikation,
  mit expliziter Trennung von Basis- und Vollständigkeitsaxiomen
  (Abschnitte~\ref{sec:axioms}, \ref{sec:fixed-point}).
\item Identifikation von (FST-7) in drei strukturell äquivalenten
  Formen (Abschnitt~\ref{sec:fst7}); die drei Formen sind in der Praxis
  drei verschiedene technische Einstiegspunkte, keine austauschbaren
  Beweisstrategien.
\item Drei technische Hindernisse (O1)--(O3) auf Route~A, katalogisiert
  mit Kandidaten-Konstruktionsskizzen, aber \emph{nicht} auf Standardtechniken
  reduziert (Abschnitt~\ref{sec:wick-mellin-obstacles}).
\item Selberg-Kompatibilitätsprüfung: Eine nicht-abelsche tannakische Erweiterung
  ist kompatibel mit Selbergs RH-Beweis
  (Abschnitt~\ref{sec:selberg}); sie re-deriviert (erzeugt) den Beweis jedoch nicht.
\item Klassifizierung von Zeta-artigen Familien
  (Tabelle~\ref{tab:classification-v2}).
\item Einen \textbf{Reduktions-Entwurf für Route~B} in der CCM-Missing-Step-Sprache
  (Entwurf~\ref{bp:fst-transfer}), zusammen mit einer
  Vier-Meilenstein-Agenda für Route~B (Meilensteine~\ref{ms:1}--\ref{ms:4}).
\item \textbf{(Neu in v0.8; durch Zookeeper vorgebracht)} Ein \textbf{Drei-Lemma-Endspiel}
  (Abschnitt~\ref{sec:three-lemma-endgame}), das (MS2/C2) im Zookeeper Paper durch
  massenbasierte spektrale Cluster-Analyse auflöst: Drei unabhängige Lemmas
  (Lemmas~\ref{lem:secular}--\ref{lem:coercive}) liefern die
  Massenkonzentrations-Schranke (Theorem~\ref{thm:mass-concentration})
  als Korollar einer Standard-Quasimode-zu-Projektor-Ungleichung.
\item \textbf{(Neu in v0.8)} Ein RFEP Transfer-Paket für physikalische
  Nachträge (Abschnitt~\ref{sec:rfep-physics-transfer}), das das
  Zookeeper-Operator-Muster in explizite Beweisverpflichtungen übersetzt:
  Operator/Form, Defekt, Galerkin-Treue, Komplement-Lücke und Rand-Limes.
\item \textbf{(Neu in v1.8)} Eine zweiglokale Status-Tabelle der Domänen
  (Tabelle~\ref{tab:domain-status}), die das RFEP als physikalische Metaregel
  vom Zeta Zoo als mathematischen Begleitzweig trennt.
\item Eine diagnostische Umformulierung von Route~C, die die
  additiv-multiplikative Kopplung (L1-K) als ihr Kernhindernis identifiziert
  (Abschnitt~\ref{sec:emergence}).
\end{enumerate}

\subsection{Was dieses Paper nicht etabliert}

\begin{enumerate}[leftmargin=*]
\item Kein unbedingtes neues Theorem über RH oder GRH.
\item Keine Ableitung von (FST-7) aus den FST-Basis-Axiomen allein;
  lediglich eine konjekturale \emph{Relokation} des arithmetischen Inputs
  unter zusätzlichen Selektionsaxiomen.
\item Keine Auflösung der Wick-Mellin-Hindernisse (O1)--(O3); die Behauptung der
  Spurklasse für (O3) aus v0.6 wird explizit zurückgezogen (Bemerkung~\ref{rem:o3-retraction}).
\item Kein unbedingter Beweis von (MS1) (Einfach-gerade für $QW_\lambda$);
  dies bleibt offen.
\item Kein unabhängiger Neu-Beweis des Zookeeper MS2/C2-Theorems innerhalb
  dieses RFEP-Papers. Dieses Paper importiert den Zookeeper-Abschluss als ein
  CORE-Resultat und nutzt ihn als Normalform-Vorlage für den physikalischen Transfer.
\item Kein automatischer Transfer des Zookeeper-Abschlusses auf physikalische
  Nachträge. Jeder physikalische Zweig muss immer noch sein eigenes Analogon
  der Hypothesen zu Operator/Form, Defekt, Quasimode, Komplement-Lücke und Domain-Limes verifizieren.
\item Kein Beweis einer generalisierten Weyl-Asymptotik für $QW_\lambda$;
  ersetzt durch Meilenstein~\ref{ms:4}.
\item Kein Beweis der physikalischen Nachträge aus dem RFEP-Transfer-Schema.
  Proposition~\ref{prop:rfep-transfer-schema} ist eine wiederverwendbare
  bedingte Vorlage; Yang--Mills, Navier--Stokes, NS-LDI, Turbulenz und CRM müssen
  ihre Hypothesen separat verifizieren.
\item Kein Beweis des nicht-abelschen Fixpunktsatzes
  (Vermutung~\ref{conj:nonabelian-fixed-point}).
\item Kein Beweis von Proposition~\ref{prop:thermo-derivation} bezüglich (B3);
  kein Beweis von Vermutung~\ref{conj:l1k} (L1-K); keine Auflösung der
  additiv-multiplikativen Kopplung.
\end{enumerate}

\subsection{Ehrliche Einschätzung}

\textbf{FST in ihrer gegenwärtigen Form ist ein vergleichendes Framework
mit einem Zookeeper-abgeschlossenen C2/MS2 Kernresultat.}
Es liefert strukturelle Klarheit, ein Klassifizierungskriterium und ein explizites
Drei-Routen-Bild, in dem der schwere offene Inhalt lokalisiert ist:

\begin{itemize}[leftmargin=*]
\item Auf Route~A sitzt der harte Inhalt in den Wick-Mellin-Hindernissen (O2), (O3).
\item Auf Route~B schließt das Zookeeper Paper das Approximationsproblem (MS2/C2)
  durch den Drei-Lemma-Mikrocluster-Abschluss. In diesem RFEP-Paper wird jenes
  Resultat importiert anstatt neu bewiesen; seine Rolle besteht darin, die Normalform für den
  physikalischen Transfer zu liefern.
\item Auf Route~C liegt der harte Inhalt in Vermutung~\ref{conj:l1k} (L1-K),
  welche das klassische Problem der additiv-multiplikativen Kopplung darstellt
  und außerhalb der Reichweite des vorliegenden Projekts liegt.
\end{itemize}

Dieses RFEP-Paper beweist die RH nicht unabhängig. Es dokumentiert
den Post-Zookeeper-Status: Das Approximationsproblem (MS2/C2) der Route~B,
welches früher als tiefer offener Punkt im CCM-Programm identifiziert wurde,
ist im Zookeeper Paper abgeschlossen. Die Konsequenz für RFEP ist strukturell:
Die physikalischen Zweige erben eine präzise Normalform, nicht jedoch eine automatische
Lösung ihrer eigenen PDE- oder Kontinuums-Limes-Probleme.

Für den physikalischen Zweig ist der Gewinn anderer Natur: RFEP hat nun eine
Transfer-Checkliste anstelle eines bloßen Slogans. Physikalische Nachträge
können das Zookeeper-Muster erst übernehmen, nachdem sie ihre Operator-Normalform,
ihren Defekt, ihren Galerkin-Limes und ihre Komplement-Lücke identifiziert haben.

\subsection{Rolle des Frameworks}

\begin{itemize}
\item \emph{Einheitliche Sprache} für Tate, Meyer, Bost-Connes, Selberg, CCM und Connes-Consani.
\item \emph{Klassifizierungskriterium} (Tabelle~\ref{tab:classification-v2}).
\item \emph{Forschungsagenda} (OFQ1--OFQ7; Meilensteine \ref{ms:1}--\ref{ms:4} für Route~B).
\item \emph{Zookeeper-Abschluss} für Route~B's (MS2/C2) durch den Drei-Lemma-Mikrocluster-Abschluss.
\item \emph{RFEP Transfer-Paket} für physikalische Nachträge (Abschnitt~\ref{sec:rfep-physics-transfer}).
\item \emph{Diagnostische Identifikation}, dass die Lücke der Route~C die additiv-multiplikative Kopplung ist.
\item \emph{Kompatibilitätsskizze} mit Selberg (Abschnitt~\ref{sec:selberg}).
\end{itemize}

\subsection{Nächste technische Schritte}

\begin{enumerate}[leftmargin=*]
\item \textbf{Route B, Meilenstein~\ref{ms:1}:} Definitions- / Formentheorie für
  $\Psi_\lambda$ als beschränkter Verknüpfungsoperator zwischen geeigneten
  Form-Definitionsbereichen. Dies ist der erste rigorose Schritt und Voraussetzung
  für jede präzise Formulierung des Entwurfs~\ref{bp:fst-transfer}.
\item \textbf{Route B, Meilenstein~\ref{ms:2}:} Min-Max-Vergleich niedriger Eigenwerte,
  keine vollständige Weyl-Asymptotik.
\item \textbf{Route B, Meilenstein~\ref{ms:3}:} Endlichdimensionales CCM-Trunkierungstheorem
  innerhalb von $E_N$.
\item \textbf{Route B, Meilenstein~\ref{ms:4}:} Erst nach 1--3 das asymptotische Regime.
\item \textbf{Physischer RFEP-Transfer:} Für jeden physikalischen Nachtrag müssen
  die fünf Beweisfelder der Proposition~\ref{prop:rfep-transfer-schema} gefüllt werden:
  Operator/Form, Defekt von endlichem Rang, Galerkin-Treue, koerzives Komplement und Rand-Limes.
  Prime-Hub ist der natürliche erste Testfall, da seine Frontier-Prime-Defekte bereits
  das richtige Endlich-Rang-/Rand-Gepräge besitzen.
\item \textbf{Route A:} Operatortheoretische Analyse von (O3) und präzise Artikulation
  der von (O2) geforderten limitierenden KMS-Zustände; die Frage nach der Spurklasse
  (Vermutung~\ref{conj:ren-trace-class}) ist offen.
\item \textbf{Route C:} Nur Dokumentation; kein weiterer Angriff ohne Fortschritte
  bei der additiv-multiplikativen Kopplung.
\item \textbf{Nicht-abelsch (OFQ3):}  Die Langlands-Funktorialität ist möglicherweise
  ein besserer langfristiger Rahmen als die gegenwärtige tannakische FST-Formulierung;
  vgl. \cite{TannakianAxioms}.
\item \textbf{Axiomatik:} Beweise für die Unabhängigkeit/Minimalität der
  Basis-Axiome; Trennung von Setup-, dynamischen und
  arithmetischen-Vollständigkeits-Annahmen (Bemerkung~\ref{rem:axioms-status}).
\end{enumerate}

Jede dieser Aufgaben ist ein mehrmonatiges Forschungsvorhaben.

\subsection{Der CRM-Kreis}

Das Kooperative Renormierungs-Modell (CRM) schließt den Kosmologie-Kreis
des RFEP v1.8. Im mathematischen Zweig ordnen eingeschränkte Positivität
und Selbstdualität die spektralen Kandidaten. Im physikalischen Zweig
beschreiben DS1--DS3, wie dissipative Dynamik einen stabilen Attraktor auswählt.
Im Kosmologie-Zweig ist das CRM die vorgeschlagene großräumige
Instanziierung: kooperative Renormierung ersetzt einen einzelnen lokalen Defekt
durch eine gekoppelte Hintergrundreaktion, und das RFEP erfordert, dass das
beobachtete Dunkle-Energie-Verhalten als stabiles Minimum des entsprechenden
Funktionals der renormierten freien Energie auftritt.

Die CRM-Behauptung ist daher absichtlich schwächer als die Behauptung
einer beobachtenden Entdeckung. Das vorliegende Paper dokumentiert
die strukturelle Verortung: Das CRM ist der kosmologische RFEP-Zweig,
dessen verbleibende Verpflichtungen die Modellidentifikation, Parameterstabilität
und beobachtbare Diskriminanten (Unterscheidungsmerkmale) sind. Diese Aufgaben
gehören in das dedizierte CRM/Kosmologie-Paper, nicht in den Beweiszweig des Zeta Zoos.

\subsection{Beziehung zum Math-Master}

Die tannakische Erweiterung des vorliegenden Papers (Abschnitt~\ref{sec:selberg})
verfeinert die HP-BL-Klassifizierung des Math-Masters \cite{MathMaster},
Abschnitte~2.4 und~9: Selbst auf der NO-Seite (Riemann, Dirichlet) könnte eine spektrale
FST-Realisierung über die SGE-YES-Route der Idele-Seite existieren. Die
Dreikomponenten-Zerlegung $\mathrm{RH}(X) \iff \mathrm{GBC}(X) \wedge C^{(2)}(X) \wedge \mathrm{Hurwitz}(X)$
des Math-Masters erhält ihre $C^{(2)}$-Komponente identifiziert mit der
Meyer-Konstruktion (wie in \cite{C2LayerStatus} diskutiert). Diese Identifikation ist
struktureller Natur und kein Beweis für die Implikationen der Zerlegung des Math-Masters.

% =====================================================
\section*{Danksagung}
% =====================================================

Dieses Paper ist ein Positionsentwurf, der im Rahmen des META\_RH\_TREE-Projekts
der FST-Mathematik erstellt wurde. Gegenüber v0.6 ist es eine ehrliche Überarbeitung,
die das externe Audit vom April 2026 (\texttt{\_audit/GPT\_REVIEW\_2026-04-17\_*.md})
einarbeitet und die Behauptungen auf Ebene des Papers in Einklang mit den
ehrlichen Diagnosen auf Ebene der Notizen bringt.
Es konsolidiert die Erkundungen von Session~11 bis Session~16, die im Projektordner
\texttt{CORE/RFEP/} dokumentiert sind, darunter:
\texttt{KONZEPT.md},
\texttt{Exploration.md},
\texttt{FST\_FIXPOINT\_THEOREM\_ATTEMPT.md},
\texttt{MINIMAL\_EXAMPLE\_RIEMANN.md},
\texttt{OPTION\_A\_PONTRJAGIN\_ABSTRACTION.md},
\texttt{OPTION\_B1\_INFORMATION\_THEORY.md},
\texttt{OPTION\_B3\_ARITHMETIC\_ERGODICITY.md},
\texttt{BRIDGE\_BOST\_CONNES\_MEYER.md},
\texttt{WICK\_ROTATION\_TECHNIQUE.md},
\texttt{O1\_DISTRIBUTION\_CLASS.md},
\texttt{O2\_FOURIER\_DUALITY.md},
\texttt{O3\_CONVERGENCE\_CONTROL.md},
\texttt{SELBERG\_CONSISTENCY\_CHECK.md},
\texttt{EMERGENCE\_HYPOTHESIS\_ATTACK.md},
\texttt{FREE\_ENERGY\_PRODUCT\_FORMULA.md},
\texttt{LEMMA\_3\_2\_PROOF.md},
\texttt{THERMODYNAMIC\_DERIVATION.md},
\texttt{TANNAKIAN\_FST\_AXIOMS.md},
\texttt{L1\_VARIATIONAL\_PROOF.md},
\texttt{L1\_K\_ATTACK.md},
\texttt{CCM\_2025\_STUDY.md},
\texttt{QW\_PW\_ISOMETRY.md} und
\texttt{ZOOKEEPER\_RFEP\_TRANSFER.md}.

% =====================================================
\begin{thebibliography}{99}

\bibitem{MathMaster}
L.~Geiger,
\textit{The Zeta Zoo --- The Mathematical Side of Functional
Stability Theory},
Preprint, Zenodo (2026), DOI: 10.5281/zenodo.19673227.

\bibitem{ZetaZooMaster}
L.~Geiger,
\textit{The Zeta Zoo --- The Mathematical Side of Functional
Stability Theory: Branch-Local Hilbert--P\'olya via the Trinity of
UCU, SGE, and Weil-Form Transversality},
Preprint, Zenodo (2026), DOI: 10.5281/zenodo.19673227.


\bibitem{RFEPTransferNote}
L.~Geiger,
\textit{Zookeeper to RFEP Transfer},
Interne Transfer-Notiz (2026),
  \texttt{CORE/RFEP/notes/ZOOKEEPER\_RFEP\_TRANSFER.md}.

\bibitem{PrimeHubSynthesis}
L.~Geiger,
\textit{Prime Hub Graph --- Synthesis of Adelic, Semiclassical,
and Bouquet Strands},
Interne Synthese-Notiz (2026),
\texttt{CORE/zetazoo/prime\_hub/SYNTHESE.md}.

\bibitem{Geiger2026RHII}
L.~Geiger,
\textit{Riemann Hypothesis via Direct Frontier Dominance
(Paper II: Even Dominance)},
Preprint (2026),
Zenodo DOI \href{https://doi.org/10.5281/zenodo.19546593}{10.5281/zenodo.19546593}.

\bibitem{AnalyticKernelV2}
L.~Geiger,
\textit{Analytic Kernel V2 --- Rigorous Weil-Formula Derivation
of the Sign Coefficient},
Ergänzende Notiz (2026),
\texttt{\_archive/\_dirichlet\_legacy/ANALYTIC\_KERNEL\_V2.md}.

\bibitem{C2LayerStatus}
L.~Geiger,
\textit{$C^{(2)}$-Layer Status Dirichlet},
Ergänzende Notiz (2026),
\texttt{CORE/zoo-mapping/C2\_LAYER\_STATUS\_2026-04-17.md}.

\bibitem{Meyer2005}
R.~Meyer,
\textit{On a representation of the idele class group related to
primes and zeros of $L$-functions},
Duke Math. J. \textbf{127} (2005), 519--595.

\bibitem{Connes1999}
A.~Connes,
\textit{Trace formula in noncommutative geometry and the zeros of
the Riemann zeta function},
Selecta Math. (N.S.) \textbf{5} (1999), 29--106.

\bibitem{Connes2026}
A.~Connes,
\textit{Recent developments in the Riemann Hypothesis programme},
arXiv:2602.04022 (2026).

\bibitem{ConnesConsaniMoscovici2025}
A.~Connes, C.~Consani, H.~Moscovici,
\textit{Zeta Spectral Triples},
arXiv:2511.22755 (2025).

\bibitem{ConnesMoscovici2022}
A.~Connes, H.~Moscovici,
\textit{The UV prolate spectrum matches the zeros of zeta},
Proc. Natl. Acad. Sci. USA \textbf{119} (22) (2022), e2123174119;
arXiv:2112.05500.

\bibitem{ConnesConsani2023}
A.~Connes, C.~Consani,
\textit{Spectral triples and $\zeta$-cycles},
Enseign. Math. \textbf{69} (2023), 93--148; arXiv:2106.01715.

\bibitem{SlepianPollack1961}
D.~Slepian, H.~O.~Pollack,
\textit{Prolate spheroidal wave functions, Fourier analysis and
uncertainty --- I},
Bell Syst. Tech. J. \textbf{40} (1961), 43--63.

\bibitem{ConnesConsaniMoscovici2024}
A.~Connes, C.~Consani, H.~Moscovici,
\textit{Zeta zeros and prolate wave operators},
arXiv:2310.18423 (2024).

\bibitem{ConnesConsaniWeilPositivity2020}
A.~Connes, C.~Consani,
\textit{Weil positivity and trace formula --- the archimedean place},
Selecta Math. (N.S.) \textbf{27} (2021), art.~77; arXiv:2006.13771.

\bibitem{Kato}
T.~Kato,
\textit{Perturbation Theory for Linear Operators},
Grundlehren der math.\ Wiss.\ \textbf{132}, Springer (1966; 2nd ed.\ 1980).

\bibitem{L1KAttack}
L.~Geiger,
\textit{L1-K --- Attack on the Reconstruction Problem},
Ergänzende Notiz (2026),
\texttt{CORE/RFEP/notes/L1\_K\_ATTACK.md}.

\bibitem{CCMStudy}
L.~Geiger,
\textit{CCM 2025 Zeta Spectral Triples --- Study + FST connection},
Ergänzende Notiz (2026),
\texttt{CORE/RFEP/notes/CCM\_2025\_STUDY.md}.

\bibitem{BostConnes1995}
J.-B.~Bost, A.~Connes,
\textit{Hecke algebras, type~III factors and phase transitions
with spontaneous symmetry breaking in number theory},
Selecta Math. (N.S.) \textbf{1} (1995), 411--457.

\bibitem{Tate1950}
J.~T.~Tate,
\textit{Fourier analysis in number fields and Hecke's
zeta-functions},
Ph.D. thesis, Princeton University (1950). Reprinted in
Cassels-Fr\"ohlich, \emph{Algebraic Number Theory}, 1967.

\bibitem{CasselsFrohlich}
J.~W.~S.~Cassels, A.~Fr\"ohlich (eds.),
\textit{Algebraic Number Theory},
Academic Press (1967).

\bibitem{ConnesMarcolli}
A.~Connes, M.~Marcolli,
\textit{Noncommutative Geometry, Quantum Fields, and Motives},
American Mathematical Society Colloquium Publications 55 (2008).

\bibitem{Selberg1956}
A.~Selberg,
\textit{Harmonic analysis and discontinuous groups in weakly
symmetric Riemannian spaces with applications to Dirichlet series},
J. Indian Math. Soc. (N.S.) \textbf{20} (1956), 47--87.

\bibitem{CornelissenMarcolli2010}
G.~Cornelissen, M.~Marcolli,
\textit{Quantum Statistical Mechanics, L-series and Anabelian
Geometry I: Partition Functions},
arXiv:1009.0736 (2010).

\bibitem{Laca1998}
M.~Laca,
\textit{Semigroups of $\ast$-endomorphisms, Dirichlet series, and
phase transitions},
J. Funct. Anal. \textbf{152} (1998), 330--378.

\bibitem{LacaRaeburn2010}
M.~Laca, I.~Raeburn,
\textit{Phase transition on the Toeplitz algebra of the affine
semigroup over the natural numbers},
Adv. Math. \textbf{225} (2010), 643--688.

\bibitem{CuntzDeningerLaca2013}
J.~Cuntz, C.~Deninger, M.~Laca,
\textit{$C^\ast$-algebras of Toeplitz type associated with
algebraic number fields},
Math. Ann. \textbf{355} (2013), 1383--1423.

\bibitem{ArtinWhaples1945}
E.~Artin, G.~Whaples,
\textit{Axiomatic characterization of fields by the product formula
for valuations},
Bull. Amer. Math. Soc. \textbf{51} (1945), 469--492.

\bibitem{EmergenceAttack}
L.~Geiger,
\textit{Emergence Hypothesis --- Attack Document},
Ergänzende Notiz (2026),
\texttt{CORE/RFEP/notes/EMERGENCE\_HYPOTHESIS\_ATTACK.md}.

\bibitem{LemmaThreeTwo}
L.~Geiger,
\textit{Lemma 3.2 --- $\N^\times$ as the unique FST$_0$-semigroup},
Ergänzende Notiz (2026),
\texttt{CORE/RFEP/notes/LEMMA\_3\_2\_PROOF.md}.

\bibitem{ThermoDerivation}
L.~Geiger,
\textit{Thermodynamic Derivation of (B1)--(B4) and FST-Minimality},
Ergänzende Notiz (2026),
\texttt{CORE/RFEP/notes/THERMODYNAMIC\_DERIVATION.md}.

\bibitem{TannakianAxioms}
L.~Geiger,
\textit{Tannakian FST Axioms --- Non-abelian Extension},
Ergänzende Notiz (2026),
\texttt{CORE/zetazoo/TANNAKIAN\_FST\_AXIOMS.md}.

\end{thebibliography}

\end{document}
